% Options for packages loaded elsewhere
\PassOptionsToPackage{unicode}{hyperref}
\PassOptionsToPackage{hyphens}{url}
\PassOptionsToPackage{dvipsnames,svgnames,x11names}{xcolor}
\PassOptionsToPackage{space}{xeCJK}
%
\documentclass[
  11pt,
]{article}
\usepackage{amsmath,amssymb}
\usepackage{iftex}
\ifPDFTeX
  \usepackage[T1]{fontenc}
  \usepackage[utf8]{inputenc}
  \usepackage{textcomp} % provide euro and other symbols
\else % if luatex or xetex
  \usepackage{unicode-math} % this also loads fontspec
  \defaultfontfeatures{Scale=MatchLowercase}
  \defaultfontfeatures[\rmfamily]{Ligatures=TeX,Scale=1}
\fi
\usepackage{lmodern}
\ifPDFTeX\else
  % xetex/luatex font selection
  \setmainfont[]{Noto Serif CJK SC}
  \setmonofont[]{Noto Sans Mono CJK SC}
  \ifXeTeX
    \usepackage{xeCJK}
    \setCJKmainfont[]{Noto Serif CJK SC}
          \fi
  \ifLuaTeX
    \usepackage[]{luatexja-fontspec}
    \setmainjfont[]{Noto Serif CJK SC}
  \fi
\fi
% Use upquote if available, for straight quotes in verbatim environments
\IfFileExists{upquote.sty}{\usepackage{upquote}}{}
\IfFileExists{microtype.sty}{% use microtype if available
  \usepackage[]{microtype}
  \UseMicrotypeSet[protrusion]{basicmath} % disable protrusion for tt fonts
}{}
\makeatletter
\@ifundefined{KOMAClassName}{% if non-KOMA class
  \IfFileExists{parskip.sty}{%
    \usepackage{parskip}
  }{% else
    \setlength{\parindent}{0pt}
    \setlength{\parskip}{6pt plus 2pt minus 1pt}}
}{% if KOMA class
  \KOMAoptions{parskip=half}}
\makeatother
\usepackage{xcolor}
\usepackage[margin=0.82in]{geometry}
\usepackage{color}
\usepackage{fancyvrb}
\newcommand{\VerbBar}{|}
\newcommand{\VERB}{\Verb[commandchars=\\\{\}]}
\DefineVerbatimEnvironment{Highlighting}{Verbatim}{commandchars=\\\{\}}
% Add ',fontsize=\small' for more characters per line
\newenvironment{Shaded}{}{}
\newcommand{\AlertTok}[1]{\textcolor[rgb]{1.00,0.00,0.00}{\textbf{#1}}}
\newcommand{\AnnotationTok}[1]{\textcolor[rgb]{0.38,0.63,0.69}{\textbf{\textit{#1}}}}
\newcommand{\AttributeTok}[1]{\textcolor[rgb]{0.49,0.56,0.16}{#1}}
\newcommand{\BaseNTok}[1]{\textcolor[rgb]{0.25,0.63,0.44}{#1}}
\newcommand{\BuiltInTok}[1]{\textcolor[rgb]{0.00,0.50,0.00}{#1}}
\newcommand{\CharTok}[1]{\textcolor[rgb]{0.25,0.44,0.63}{#1}}
\newcommand{\CommentTok}[1]{\textcolor[rgb]{0.38,0.63,0.69}{\textit{#1}}}
\newcommand{\CommentVarTok}[1]{\textcolor[rgb]{0.38,0.63,0.69}{\textbf{\textit{#1}}}}
\newcommand{\ConstantTok}[1]{\textcolor[rgb]{0.53,0.00,0.00}{#1}}
\newcommand{\ControlFlowTok}[1]{\textcolor[rgb]{0.00,0.44,0.13}{\textbf{#1}}}
\newcommand{\DataTypeTok}[1]{\textcolor[rgb]{0.56,0.13,0.00}{#1}}
\newcommand{\DecValTok}[1]{\textcolor[rgb]{0.25,0.63,0.44}{#1}}
\newcommand{\DocumentationTok}[1]{\textcolor[rgb]{0.73,0.13,0.13}{\textit{#1}}}
\newcommand{\ErrorTok}[1]{\textcolor[rgb]{1.00,0.00,0.00}{\textbf{#1}}}
\newcommand{\ExtensionTok}[1]{#1}
\newcommand{\FloatTok}[1]{\textcolor[rgb]{0.25,0.63,0.44}{#1}}
\newcommand{\FunctionTok}[1]{\textcolor[rgb]{0.02,0.16,0.49}{#1}}
\newcommand{\ImportTok}[1]{\textcolor[rgb]{0.00,0.50,0.00}{\textbf{#1}}}
\newcommand{\InformationTok}[1]{\textcolor[rgb]{0.38,0.63,0.69}{\textbf{\textit{#1}}}}
\newcommand{\KeywordTok}[1]{\textcolor[rgb]{0.00,0.44,0.13}{\textbf{#1}}}
\newcommand{\NormalTok}[1]{#1}
\newcommand{\OperatorTok}[1]{\textcolor[rgb]{0.40,0.40,0.40}{#1}}
\newcommand{\OtherTok}[1]{\textcolor[rgb]{0.00,0.44,0.13}{#1}}
\newcommand{\PreprocessorTok}[1]{\textcolor[rgb]{0.74,0.48,0.00}{#1}}
\newcommand{\RegionMarkerTok}[1]{#1}
\newcommand{\SpecialCharTok}[1]{\textcolor[rgb]{0.25,0.44,0.63}{#1}}
\newcommand{\SpecialStringTok}[1]{\textcolor[rgb]{0.73,0.40,0.53}{#1}}
\newcommand{\StringTok}[1]{\textcolor[rgb]{0.25,0.44,0.63}{#1}}
\newcommand{\VariableTok}[1]{\textcolor[rgb]{0.10,0.09,0.49}{#1}}
\newcommand{\VerbatimStringTok}[1]{\textcolor[rgb]{0.25,0.44,0.63}{#1}}
\newcommand{\WarningTok}[1]{\textcolor[rgb]{0.38,0.63,0.69}{\textbf{\textit{#1}}}}
\usepackage{longtable,booktabs,array}
\usepackage{calc} % for calculating minipage widths
% Correct order of tables after \paragraph or \subparagraph
\usepackage{etoolbox}
\makeatletter
\patchcmd\longtable{\par}{\if@noskipsec\mbox{}\fi\par}{}{}
\makeatother
% Allow footnotes in longtable head/foot
\IfFileExists{footnotehyper.sty}{\usepackage{footnotehyper}}{\usepackage{footnote}}
\makesavenoteenv{longtable}
\setlength{\emergencystretch}{3em} % prevent overfull lines
\providecommand{\tightlist}{%
  \setlength{\itemsep}{0pt}\setlength{\parskip}{0pt}}
\setcounter{secnumdepth}{-\maxdimen} % remove section numbering
\ifLuaTeX
  \usepackage{selnolig}  % disable illegal ligatures
\fi
\usepackage{bookmark}
\IfFileExists{xurl.sty}{\usepackage{xurl}}{} % add URL line breaks if available
\urlstyle{same}
\hypersetup{
  colorlinks=true,
  linkcolor={Maroon},
  filecolor={Maroon},
  citecolor={Blue},
  urlcolor={Blue},
  pdfcreator={LaTeX via pandoc}}

\author{}
\date{}

\begin{document}

\section{MATH1090
期中考试详细复习讲义}\label{math1090-ux671fux4e2dux8003ux8bd5ux8be6ux7ec6ux590dux4e60ux8bb2ux4e49}

\textbf{课程}：MATH1090 / 1098 (Honours) Introduction to Set Theory\\
\textbf{版本}：Midterm Review Edition\\
\textbf{范围}：Lecture Notes 至 §4.7；\textbf{Dedekind cuts
不纳入期中}\\
\textbf{语言说明}：全文以简体中文为主；专有名词首次出现时附英文\\
\textbf{主要材料}：Lecture Notes, Homework 1, Homework 2, Worksheet 1--5

\subsection{文档导航}\label{ux6587ux6863ux5bfcux822a}

\begin{itemize}
\tightlist
\item
  第一章：逻辑（命题逻辑、谓词逻辑、HW1、HW2、Worksheet 1--2）
\item
  第二章：集合、函数与关系（Lecture Exercises 11--24，Worksheet 3）
\item
  第三章：自然数、整数、有理数（Lecture Exercises 25--38，Worksheet
  4--5）
\item
  第四章：实数导论至 §4.7（Lecture Exercises 39--44，Q 不完备，完备性）
\item
  第五章：高频证明模板与临考整理
\item
  第六章：最后一遍必看清单
\item
  附录：定义速查、快问快答、半小时速记版、补充严解
\end{itemize}

\begin{center}\rule{0.5\linewidth}{0.5pt}\end{center}

\subsection{使用说明}\label{ux4f7fux7528ux8bf4ux660e}

这份讲义的目标不是把 lecture notes
简单翻译一遍，而是把期中之前的所有核心概念、命题、常见证明套路、以及已提供作业与
worksheet 的题目全部整理到一份可以直接复习与背诵的文档里。\\
我采用的组织方式是：

\begin{enumerate}
\def\labelenumi{\arabic{enumi}.}
\tightlist
\item
  先给每一章的\textbf{概念总览}，帮助你建立结构；
\item
  再给\textbf{必须掌握的定义、命题与证明思路}；
\item
  最后把\textbf{lecture notes / homework / worksheet}
  里的题目逐题写出详细解答。
\end{enumerate}

你复习时建议按下面顺序进行：

\begin{itemize}
\tightlist
\item
  第一遍：先看每章前面的``核心总结''，把概念网络串起来；
\item
  第二遍：专门看``证明模板''，熟悉如何写严谨证明；
\item
  第三遍：遮住解答，自己做题，再对照答案查漏补缺；
\item
  最后一天：只看``高频结论清单''和自己最容易混淆的地方。
\end{itemize}

\subsection{期中范围地图（Course
Map）}\label{ux671fux4e2dux8303ux56f4ux5730ux56fecourse-map}

\begin{longtable}[]{@{}
  >{\raggedright\arraybackslash}p{(\columnwidth - 4\tabcolsep) * \real{0.3333}}
  >{\raggedright\arraybackslash}p{(\columnwidth - 4\tabcolsep) * \real{0.3333}}
  >{\raggedright\arraybackslash}p{(\columnwidth - 4\tabcolsep) * \real{0.3333}}@{}}
\toprule\noalign{}
\begin{minipage}[b]{\linewidth}\raggedright
模块
\end{minipage} & \begin{minipage}[b]{\linewidth}\raggedright
范围
\end{minipage} & \begin{minipage}[b]{\linewidth}\raggedright
你要会什么
\end{minipage} \\
\midrule\noalign{}
\endhead
\bottomrule\noalign{}
\endlastfoot
逻辑 (Logic) & Chapter 1 &
命题逻辑、谓词逻辑、量词、否定、逻辑等价、推理有效性 \\
集合 (Sets) & Chapter 2 前半 &
集合运算、子集、幂集、笛卡尔积、集合代数 \\
函数与关系 (Functions and Relations) & Chapter 2 后半 &
函数、单射满射双射、逆、关系、偏序、等价关系、商集 \\
自然数与整数 (Naturals and Integers) & Chapter 3 前半 & Peano
公理、递归定义加法乘法、归纳法、从 N 构造 Z \\
有理数 (Rationals) & Chapter 3 后半 & 从 Z 构造 Q、运算良定义、逆元、√2
非有理 \\
实数导论 (The Reals) & Chapter 4 至 §4.7 &
全序、上界下界、sup/inf、完备性、Q 不完备、实数公理化、4.7
的``第一次尝试'' \\
\end{longtable}

\subsection{常用符号表}\label{ux5e38ux7528ux7b26ux53f7ux8868}

\begin{longtable}[]{@{}lll@{}}
\toprule\noalign{}
符号 & 名称 & 含义 \\
\midrule\noalign{}
\endhead
\bottomrule\noalign{}
\endlastfoot
¬A & 否定 (negation) & A 不真 \\
A ∧ B & 合取 (conjunction) & A 且 B \\
A ∨ B & 析取 (disjunction) & A 或 B（包含至少一个为真） \\
A → B & 蕴含 (implication) & 若 A 则 B \\
A ↔ B & 双条件 (biconditional) & A 当且仅当 B \\
∀x & 全称量词 (universal quantifier) & 对所有 x \\
∃x & 存在量词 (existential quantifier) & 存在某个 x \\
∃!x & 唯一存在 (exists uniquely) & 存在且仅存在一个 x \\
A ⊂ B & 子集 (subset) & A 中每个元素都在 B 中 \\
A ∪ B & 并集 (union) & 在 A 或在 B \\
A ∩ B & 交集 (intersection) & 同时在 A 与 B \\
\(A\setminus B\) & 差集 (set difference) & 在 A 但不在 B \\
A × B & 笛卡尔积 (Cartesian product) & 有序对集合 \\
P(A) & 幂集 (power set) & A 的所有子集 \\
f : X → Y & 函数 (function) & 从 X 到 Y 的映射 \\
\(id_X\) & 恒等映射 (identity) & \(x \mapsto x\) \\
\(\preceq\) & 偏序关系符号 & partial order 的写法 \\
{[}x{]}\_R & 等价类 (equivalence class) & 与 x 等价的所有元素 \\
X / \textasciitilde{} & 商集 (quotient set) & 等价类构成的集合 \\
sup(Y) & 上确界 (supremum) & 最小上界 \\
inf(Y) & 下确界 (infimum) & 最大下界 \\
\end{longtable}

\begin{center}\rule{0.5\linewidth}{0.5pt}\end{center}

\section{第一章
逻辑（Logic）}\label{ux7b2cux4e00ux7ae0-ux903bux8f91logic}

\subsection{1.1 核心总结}\label{ux6838ux5fc3ux603bux7ed3}

这一章的核心任务有两件事：

\begin{itemize}
\tightlist
\item
  把自然语言论证翻译成形式语言；
\item
  判断一个公式在什么情况下为真，以及某个结论是否能从前提推出。
\end{itemize}

命题逻辑 (propositional logic) 研究的是``整句话''的真假；谓词逻辑
(predicate logic)
则进一步研究``依赖于变量的性质''，因此能够表达``对所有''``存在''等结构。

\subsubsection{本章必背结论}\label{ux672cux7ae0ux5fc5ux80ccux7ed3ux8bba}

\begin{enumerate}
\def\labelenumi{\arabic{enumi}.}
\tightlist
\item
  运算优先级：先 ¬，再 ∧ / ∨，最后 → / ↔。\\
\item
  蕴含可以改写为：\\
  A → B ≡ ¬A ∨ B
\item
  逆否命题等价：\\
  A → B ≡ ¬B → ¬A
\item
  德摩根律 (de Morgan's laws)：\\
  ¬(A ∧ B) ≡ ¬A ∨ ¬B\\
  ¬(A ∨ B) ≡ ¬A ∧ ¬B
\item
  量词否定：\\
  ¬∀x P(x) ≡ ∃x ¬P(x)\\
  ¬∃x P(x) ≡ ∀x ¬P(x)
\item
  混合量词一般\textbf{不能换序}：\\
  ∀x∃y P(x,y) 与 ∃y∀x P(x,y) 通常不等价。
\end{enumerate}

\subsection{1.2 命题逻辑（Propositional
Logic）概念整理}\label{ux547dux9898ux903bux8f91propositional-logicux6982ux5ff5ux6574ux7406}

\subsubsection{命题（proposition）}\label{ux547dux9898proposition}

命题是一个\textbf{有确定真假值}的陈述句。\\
例如：

\begin{itemize}
\tightlist
\item
  ``2+2=4'' 是命题，而且为真；
\item
  ``关门！'' 不是命题，因为它不是陈述真假；
\item
  ``x+1=3'' 若没有说明 x 是什么，一般不视为命题，而更像是开放句 (open
  sentence)。
\end{itemize}

\subsubsection{五个基本联结词（connectives）}\label{ux4e94ux4e2aux57faux672cux8054ux7ed3ux8bcdconnectives}

\begin{longtable}[]{@{}lll@{}}
\toprule\noalign{}
记号 & 英文 & 中文理解 \\
\midrule\noalign{}
\endhead
\bottomrule\noalign{}
\endlastfoot
¬A & NOT & 非 A \\
A ∧ B & AND & A 且 B \\
A ∨ B & OR & A 或 B \\
A → B & IMPLIES & 若 A 则 B \\
A ↔ B & IFF & A 当且仅当 B \\
\end{longtable}

\subsubsection{对蕴含 A → B
的理解}\label{ux5bf9ux8574ux542b-a-b-ux7684ux7406ux89e3}

A → B 只有在 \textbf{A 真而 B 假} 时为假，其余都为真。\\
所以它表达的不是`` A 已经发生''，而是``只要 A 发生，B 就必须发生''。

这也是初学时最容易错的地方：\\
``如果下雨，则地面湿''并\textbf{不}表示现在真的在下雨。

\subsection{1.3 真值表（Truth
Tables）总复习}\label{ux771fux503cux8868truth-tablesux603bux590dux4e60}

\subsubsection{单个公式的真值表}\label{ux5355ux4e2aux516cux5f0fux7684ux771fux503cux8868}

\begin{longtable}[]{@{}ll@{}}
\toprule\noalign{}
A & ¬A \\
\midrule\noalign{}
\endhead
\bottomrule\noalign{}
\endlastfoot
F & T \\
T & F \\
\end{longtable}

\begin{longtable}[]{@{}llll@{}}
\toprule\noalign{}
A & B & A ∧ B & A ∨ B \\
\midrule\noalign{}
\endhead
\bottomrule\noalign{}
\endlastfoot
F & F & F & F \\
F & T & F & T \\
T & F & F & T \\
T & T & T & T \\
\end{longtable}

\begin{longtable}[]{@{}llll@{}}
\toprule\noalign{}
A & B & A → B & A ↔ B \\
\midrule\noalign{}
\endhead
\bottomrule\noalign{}
\endlastfoot
F & F & T & T \\
F & T & T & F \\
T & F & F & F \\
T & T & T & T \\
\end{longtable}

\subsubsection{如何判断逻辑等价（logical
equivalence）}\label{ux5982ux4f55ux5224ux65adux903bux8f91ux7b49ux4ef7logical-equivalence}

两个公式逻辑等价，意思是：\\
\textbf{无论命题变量取什么真假值，它们的最终真值都相同。}

判断方法只有一个最扎实的：\textbf{列真值表}。\\
当然熟练之后也可以直接用已知等价式推导。

\subsection{1.4
命题逻辑常见等价式}\label{ux547dux9898ux903bux8f91ux5e38ux89c1ux7b49ux4ef7ux5f0f}

\begin{longtable}[]{@{}
  >{\raggedright\arraybackslash}p{(\columnwidth - 2\tabcolsep) * \real{0.5000}}
  >{\raggedright\arraybackslash}p{(\columnwidth - 2\tabcolsep) * \real{0.5000}}@{}}
\toprule\noalign{}
\begin{minipage}[b]{\linewidth}\raggedright
名称
\end{minipage} & \begin{minipage}[b]{\linewidth}\raggedright
公式
\end{minipage} \\
\midrule\noalign{}
\endhead
\bottomrule\noalign{}
\endlastfoot
双重否定 & ¬(¬A) ≡ A \\
蕴含改写 & A → B ≡ ¬A ∨ B \\
交换律 & A ∧ B ≡ B ∧ A；A ∨ B ≡ B ∨ A \\
结合律 & A ∧ (B ∧ C) ≡ (A ∧ B) ∧ C；A ∨ (B ∨ C) ≡ (A ∨ B) ∨ C \\
分配律 & A ∧ (B ∨ C) ≡ (A ∧ B) ∨ (A ∧ C)；A ∨ (B ∧ C) ≡ (A ∨ B) ∧ (A ∨
C) \\
德摩根律 & ¬(A ∧ B) ≡ ¬A ∨ ¬B；¬(A ∨ B) ≡ ¬A ∧ ¬B \\
逆否等价 & A → B ≡ ¬B → ¬A \\
\end{longtable}

\subsubsection{tautology 与
contradiction}\label{tautology-ux4e0e-contradiction}

\begin{itemize}
\tightlist
\item
  永真式 (tautology)：无论怎样代入都为真。\\
  例如 A ∨ ¬A。
\item
  永假式 (contradiction)：无论怎样代入都为假。\\
  例如 A ∧ ¬A。
\end{itemize}

\subsection{1.5 命题逻辑推理（Rules of
Inference）与常见谬误（Fallacies）}\label{ux547dux9898ux903bux8f91ux63a8ux7406rules-of-inferenceux4e0eux5e38ux89c1ux8c2cux8beffallacies}

\subsubsection{常见有效推理}\label{ux5e38ux89c1ux6709ux6548ux63a8ux7406}

\begin{enumerate}
\def\labelenumi{\arabic{enumi}.}
\item
  \textbf{假言三段论}\\
  A → B\\
  B → C\\
  ∴ A → C
\item
  \textbf{否后式 / modus tollens} B → A\\
  ¬A\\
  ∴ ¬B
\item
  \textbf{析取三段论} A ∨ B\\
  ¬A\\
  ∴ B
\item
  \textbf{去合取消去} (A ∨ B) ∧ ¬B\\
  ∴ A
\end{enumerate}

\subsubsection{常见错误推理}\label{ux5e38ux89c1ux9519ux8befux63a8ux7406}

\begin{enumerate}
\def\labelenumi{\arabic{enumi}.}
\item
  \textbf{肯后式（affirming the consequent）} A → B\\
  B\\
  ∴ A\\
  错，因为 B 也可能由别的原因导致。
\item
  \textbf{否前式（denying the antecedent）} A → B\\
  ¬A\\
  ∴ ¬B\\
  错，因为 B 仍可能独立为真。
\end{enumerate}

\begin{center}\rule{0.5\linewidth}{0.5pt}\end{center}

\subsection{1.6 Lecture Notes 习题详解（Exercises
1--10）}\label{lecture-notes-ux4e60ux9898ux8be6ux89e3exercises-110}

\subsubsection{Exercise 1}\label{exercise-1}

\textbf{题意}：就 murder 案例中的 A, B, C, D, E, F, G, H
再写出一些逻辑依赖关系。

\textbf{解答}：这是开放题，关键是写出``合理的布尔公式 (Boolean
formulas)''。例如：

\begin{itemize}
\tightlist
\item
  A ↔ ¬B\\
  表示``案发在夜间''与``案发不在白天''等价。
\item
  F → A\\
  表示``如果 Omar 是凶手，那么案发在夜间''，这是因为已知 C 说 Omar
  夜里在睡觉；如果进一步结合``凶手作案时不能在睡觉''，可以建立这样的依赖。
\item
  G → B\\
  类似地，如果 Michael 是凶手，则更可能与白天案发有关。
\item
  A ∧ C → ¬F\\
  若案发在夜间且 Omar 夜里在睡觉，则 Omar 不可能是凶手。
\item
  B ∧ D → ¬G\\
  这是 notes 已经给出的公式。
\item
  E ∧ A → ¬H\\
  若案发时刻落在 Mykola 睡觉时段，则 Mykola 不是凶手。
\end{itemize}

\textbf{点评}：本题没有唯一答案，重点不在``猜故事''，而在于把自然语言条件翻译成符号。

\subsubsection{Exercise 2}\label{exercise-2}

\textbf{题意}：验证 notes 中列出的逻辑等价式。

\textbf{解答}：做法是分别对左右两边列真值表，比较最后一列是否完全相同。这里给出最重要的几个验证结论：

\begin{enumerate}
\def\labelenumi{\arabic{enumi}.}
\item
  \textbf{蕴含改写} A → B 与 ¬A ∨ B 的真值在四种情形下完全一致，所以\\
  A → B ≡ ¬A ∨ B。
\item
  \textbf{德摩根律} ¬(A ∧ B) 与 ¬A ∨ ¬B 的真值完全一致；\\
  ¬(A ∨ B) 与 ¬A ∧ ¬B 的真值完全一致。
\item
  \textbf{逆否命题} A → B 与 ¬B → ¬A 的真值表也完全相同。
\item
  \textbf{交换律、结合律、分配律} 这些都可通过 2 变量或 3
  变量真值表验证。
\end{enumerate}

\textbf{复习建议}：期中不一定要求你每次都从头列完整表，但至少要会对最核心的三个：\\
A → B ≡ ¬A ∨ B，\\
德摩根律，\\
逆否等价。

\subsubsection{Exercise 3}\label{exercise-3}

\textbf{题意}：用真值表检查前面两个推理例子。

\textbf{解答}：

\paragraph{例 1}\label{ux4f8b-1}

前提： - A → B - B → C

结论： - A → C

若前提都真，则不可能出现 A 真且 C 假。\\
因为若 A 真，由 A → B 得 B 真；再由 B → C 得 C 真。\\
所以 A → C 必真。故推理有效。

\paragraph{例 2}\label{ux4f8b-2}

前提： - ¬A - B → A

结论： - ¬B

如果 B 真，则由 B → A 得 A 真；但前提又给 ¬A。矛盾。\\
所以 B 不真，即 ¬B 真。故推理有效。

\subsubsection{Exercise 4}\label{exercise-4}

\textbf{题意}：若 (A ∨ B) ∧ ¬B 为真，推出 A。

\textbf{解答}：\\
由 (A ∨ B) ∧ ¬B 为真，可知：

\begin{itemize}
\tightlist
\item
  A ∨ B 为真；
\item
  ¬B 为真，因此 B 为假。
\end{itemize}

既然 A ∨ B 为真，而 B 为假，则只能 A 为真。\\
所以结论 A 成立。\\
这正是``析取消去''或``析取三段论''的经典形式。

\subsubsection{Exercise 5}\label{exercise-5}

\textbf{题意}：证明\\
A → B, B → C, C → A\\
与 ``A ↔ B ↔ C'' 等价。

\textbf{解答}：\\
在本课程语境下，``A ↔ B ↔ C'' 可理解为``这三个命题真假相同''，等价于\\
(A ↔ B) ∧ (B ↔ C)。

而 - A ↔ B 等价于 (A → B) ∧ (B → A) - B ↔ C 等价于 (B → C) ∧ (C → B)

如果我们已知 A → B, B → C, C → A， 则由传递性可得： - B → A（因为 B → C
且 C → A） - C → B（因为 C → A 且 A → B）

于是得到 A ↔ B 且 B ↔ C。\\
反过来，如果 A ↔ B 且 B ↔ C，则显然有 A → B、B → C、C → A。\\
故两者等价。

\subsubsection{Exercise 6}\label{exercise-6}

\textbf{题意}：把下列句子写成谓词公式。

\textbf{解答}：先选定合适的论域 (domain of discourse)。

\begin{enumerate}
\def\labelenumi{\arabic{enumi}.}
\item
  ``所有情侣都会吵架。''\\
  设 C(x,y)：x 与 y 是情侣； F(x,y)：x 与 y 会吵架。\\
  则可写为\\
  ∀x∀y (C(x,y) → F(x,y))。
\item
  ``每个人至少在读两本书。''\\
  设 R(x,b)：x 在读书 b。\\
  则可写为\\
  \(\forall x\, \exists b_1\, \exists b_2\, (b_1 \neq b_2 \land R(x,b_1) \land R(x,b_2))\)。
\item
  ``我的同学恰好有一个朋友。''\\
  设 o 表示``我的同学''，F(o,x) 表示 x 是 o 的朋友。\\
  可写为\\
  ∃x (F(o,x) ∧ ∀y (F(o,y) → y = x))。\\
  也可以写成\\
  ∃!x F(o,x)。
\end{enumerate}

\subsubsection{Exercise 7}\label{exercise-7}

\textbf{题意}：把 ¬∀xP(x) 用 ∃ 改写。

\textbf{解答}：\\
¬∀xP(x) ≡ ∃x¬P(x)。

\textbf{解释}：\\
``不是所有 x 都满足 P'' 等价于``至少存在一个 x 不满足 P''。

\subsubsection{Exercise 8}\label{exercise-8}

\textbf{题意}：写出一些谓词逻辑中的``规律''（laws）。

\textbf{解答}：常见规律包括：

\begin{itemize}
\tightlist
\item
  量词否定：

  \begin{itemize}
  \tightlist
  \item
    ¬∀xP(x) ≡ ∃x¬P(x)
  \item
    ¬∃xP(x) ≡ ∀x¬P(x)
  \end{itemize}
\item
  同类量词换序：

  \begin{itemize}
  \tightlist
  \item
    ∀x∀yP(x,y) ≡ ∀y∀xP(x,y)
  \item
    ∃x∃yP(x,y) ≡ ∃y∃xP(x,y)
  \end{itemize}
\item
  分配型公式：

  \begin{itemize}
  \tightlist
  \item
    ∀x(P(x) ∧ Q(x)) ≡ (∀xP(x)) ∧ (∀xQ(x))
  \item
    ∃x(P(x) ∨ Q(x)) ≡ (∃xP(x)) ∨ (∃xQ(x))
  \end{itemize}
\end{itemize}

\textbf{注意}：不要把``看起来很像''的公式都当作恒等式，混合量词和 ∨/∧
的组合最容易出错。

\subsubsection{Exercise 9}\label{exercise-9}

\textbf{题意}：翻译并比较\\
∀x∃y(x\textless y) 与 ∃y∀x(x\textless y)。

\textbf{解答}：

\begin{itemize}
\item
  ∀x∃y(x \textless{} y)：\\
  ``对每一个 x，都存在一个比它更大的 y。''\\
  这表示``这个集合中没有最大元''。
\item
  ∃y∀x(x \textless{} y)：\\
  ``存在一个 y，使得所有 x 都小于 y。''\\
  这表示``存在一个统一的最大上方元素''。
\end{itemize}

二者差别在于：\\
第一个句子里 y \textbf{依赖于} x；\\
第二个句子里 y 是\textbf{同一个固定对象}，要对全部 x 一起工作。\\
因此它们通常不等价。

\subsubsection{Exercise 10}\label{exercise-10}

\textbf{题意}：列出一些``看起来像对，但其实失败''的类似公式。

\textbf{解答}：下面是几个非常经典的失败例子：

\begin{enumerate}
\def\labelenumi{\arabic{enumi}.}
\item
  ∀x(P(x) ∨ Q(x)) \textbf{不一定等价于} (∀xP(x)) ∨ (∀xQ(x))。\\
  反例：有人喜欢数学，有人喜欢英语，但不是所有人都喜欢数学，也不是所有人都喜欢英语。
\item
  ∃x(P(x) ∧ Q(x)) \textbf{不一定等价于} (∃xP(x)) ∧ (∃xQ(x))。\\
  右边只要求``存在某个 P 的见证''和``存在某个 Q
  的见证''，这两个见证可以不是同一个人。
\item
  ∀x∃yP(x,y) \textbf{不一定等价于} ∃y∀xP(x,y)。\\
  这就是 Exercise 9 的本质。
\end{enumerate}

\begin{center}\rule{0.5\linewidth}{0.5pt}\end{center}

\subsection{1.7 Homework 1 详解（Propositional
Logic）}\label{homework-1-ux8be6ux89e3propositional-logic}

\subsubsection{HW1 第 1
题：补括号并翻译}\label{hw1-ux7b2c-1-ux9898ux8865ux62ecux53f7ux5e76ux7ffbux8bd1}

\paragraph{(a) ¬A ∧ B ∨ C}\label{a-a-b-c}

按优先级应读作： (¬A ∧ B) ∨ C

中文：\\
``A 不真且 B 为真，或者 C 为真。''

\paragraph{(b) A → B ∧ ¬C}\label{b-a-b-c}

按优先级应读作： A → (B ∧ ¬C)

中文：\\
``如果 A 成立，那么 B 成立且 C 不成立。''

\paragraph{(c) ¬A ↔ B ∨ C}\label{c-a-b-c}

按优先级应读作： ¬A ↔ (B ∨ C)

中文：\\
``A 不成立，当且仅当 B 或 C 至少一个成立。''

\subsubsection{HW1 第 2
题：真值表}\label{hw1-ux7b2c-2-ux9898ux771fux503cux8868}

\paragraph{(a) (A ∨ B) → C}\label{a-a-b-c-1}

\begin{longtable}[]{@{}lllll@{}}
\toprule\noalign{}
A & B & C & A ∨ B & (A ∨ B) → C \\
\midrule\noalign{}
\endhead
\bottomrule\noalign{}
\endlastfoot
F & F & F & F & T \\
F & F & T & F & T \\
F & T & F & T & F \\
F & T & T & T & T \\
T & F & F & T & F \\
T & F & T & T & T \\
T & T & F & T & F \\
T & T & T & T & T \\
\end{longtable}

\paragraph{(b) A → (B → A)}\label{b-a-b-a}

\begin{longtable}[]{@{}llll@{}}
\toprule\noalign{}
A & B & B → A & A → (B → A) \\
\midrule\noalign{}
\endhead
\bottomrule\noalign{}
\endlastfoot
F & F & T & T \\
F & T & F & T \\
T & F & T & T \\
T & T & T & T \\
\end{longtable}

所以它是永真式。

\paragraph{(c) (A ∧ B) ↔ (B ∧ A)}\label{c-a-b-b-a}

由于 ∧ 满足交换律，左右恒相同。\\
真值表最后一列恒为 T，因此也是永真式。

\subsubsection{HW1 第 3 题：判断 tautology / contradiction /
neither}\label{hw1-ux7b2c-3-ux9898ux5224ux65ad-tautology-contradiction-neither}

\begin{enumerate}
\def\labelenumi{\arabic{enumi}.}
\tightlist
\item
  A ∨ ¬A：永真式 (tautology)。
\item
  A ∧ ¬A：永假式 (contradiction)。
\item
  (A → B) ∨ (B → A)：永真式。\\
  因为如果 A 真 B 假，则 B → A 为真；\\
  如果 A 假 B 真，则 A → B 为真；\\
  其余情况两个都真。
\end{enumerate}

\subsubsection{HW1 第 4
题：证明逻辑等价}\label{hw1-ux7b2c-4-ux9898ux8bc1ux660eux903bux8f91ux7b49ux4ef7}

\begin{enumerate}
\def\labelenumi{\arabic{enumi}.}
\item
  ¬(A ∨ B) ≡ ¬A ∧ ¬B\\
  这是德摩根律。
\item
  A ↔ B ≡ ¬A ↔ ¬B\\
  因为 A 与 B 同真同假，当且仅当 ¬A 与 ¬B 也同真同假。
\end{enumerate}

\subsubsection{HW1 第 5
题：去掉蕴含符号}\label{hw1-ux7b2c-5-ux9898ux53bbux6389ux8574ux542bux7b26ux53f7}

¬(A → B) ≡ ¬(¬A ∨ B) ≡ A ∧ ¬B

这是非常重要的一条：\\
``否定一个蕴含''等价于``前件真且后件假''。

\subsubsection{HW1 第 6 题：判断 deductions
是否有效}\label{hw1-ux7b2c-6-ux9898ux5224ux65ad-deductions-ux662fux5426ux6709ux6548}

\paragraph{(a)}\label{a}

A → B\\
B → C\\
∴ A → C

有效。假言三段论。

\paragraph{(b)}\label{b}

A ∨ B\\
¬A\\
∴ B

有效。析取三段论。

\paragraph{(c)}\label{c}

A → B\\
¬A\\
∴ ¬B

无效。否前式 (denying the antecedent)。\\
反例：\\
A = ``我在香港''，B = ``我在地球上''。\\
A → B 为真，¬A 也真，但 ¬B 假。

\begin{center}\rule{0.5\linewidth}{0.5pt}\end{center}

\subsection{1.8 Worksheet 1 详解（Propositional Logic in
Practice）}\label{worksheet-1-ux8be6ux89e3propositional-logic-in-practice}

\subsubsection{Exercise 1：判断是不是
proposition}\label{exercise-1ux5224ux65adux662fux4e0dux662f-proposition}

\begin{enumerate}
\def\labelenumi{\arabic{enumi}.}
\tightlist
\item
  2 + 2 = 4\\
  是命题，且为真。
\item
  x + 1 = 3\\
  若 x 未指定，不是命题。
\item
  7 是素数\\
  是命题，且为真。
\item
  Close the door.\\
  不是命题，是命令句。
\item
  Every even number is divisible by 2.\\
  是命题，且为真。
\end{enumerate}

\subsubsection{Exercise 2：翻译
connectives}\label{exercise-2ux7ffbux8bd1-connectives}

\begin{itemize}
\tightlist
\item
  P ∧ Q：P 且 Q。
\item
  P ∨ Q：P 或 Q。
\item
  ¬P：非 P。
\item
  P → Q：如果 P，则 Q。
\item
  P ↔ Q：P 当且仅当 Q。
\end{itemize}

\subsubsection{Exercise 3：English to
Symbols}\label{exercise-3english-to-symbols}

设\\
P = ``作业已提交''\\
Q = ``作业已批改''

\begin{enumerate}
\def\labelenumi{\arabic{enumi}.}
\tightlist
\item
  已提交且已批改：P ∧ Q\\
\item
  已提交但未批改：P ∧ ¬Q\\
\item
  如果已批改，则它曾被提交：Q → P\\
\item
  已批改当且仅当已提交：Q ↔ P
\end{enumerate}

\subsubsection{Exercise 4：真值表}\label{exercise-4ux771fux503cux8868}

\paragraph{(a) ¬P ∨ Q}\label{a-p-q}

这与 P → Q 等价。

\begin{longtable}[]{@{}llll@{}}
\toprule\noalign{}
P & Q & ¬P & ¬P ∨ Q \\
\midrule\noalign{}
\endhead
\bottomrule\noalign{}
\endlastfoot
F & F & T & T \\
F & T & T & T \\
T & F & F & F \\
T & T & F & T \\
\end{longtable}

\paragraph{(b) P → Q}\label{b-p-q}

与上表最后一列相同。

\paragraph{(c) (P ∧ Q) → P}\label{c-p-q-p}

只要前件真，P 必真；若前件假，蕴含也真。\\
所以它是永真式。

\paragraph{(d) P → (Q → P)}\label{d-p-q-p}

不管 Q 怎样，只要 P 真，则 Q → P 真；若 P 假，则外层蕴含自动真。\\
因此也是永真式。

\subsubsection{Exercise 5：tautology / contradiction /
neither}\label{exercise-5tautology-contradiction-neither}

\begin{itemize}
\tightlist
\item
  P ∨ ¬P：tautology
\item
  P ∧ ¬P：contradiction
\item
  (P ∧ Q) → P：tautology
\item
  P → Q：neither
\end{itemize}

\subsubsection{Exercise 6：理解
implication}\label{exercise-6ux7406ux89e3-implication}

设\\
P = ``正在下雨''\\
Q = ``地面是湿的''

\begin{enumerate}
\def\labelenumi{\arabic{enumi}.}
\tightlist
\item
  P → Q：如果下雨，则地面湿。
\item
  它并不表示``现在正在下雨''，因为当前件假时整个蕴含仍可能为真。
\item
  例子：今天没下雨，但洒水车刚开过，所以地面湿。\\
  此时 P 假，P → Q 仍真。
\end{enumerate}

\subsubsection{Exercise
7：常见逻辑错误}\label{exercise-7ux5e38ux89c1ux903bux8f91ux9519ux8bef}

\paragraph{(a) If P then Q. Q. Therefore,
P.}\label{a-if-p-then-q.-q.-therefore-p.}

这是\textbf{肯后式} (affirming the consequent)。\\
反例：\\
如果我是数学系学生，则我是学生。\\
我是学生。\\
所以我是数学系学生。\\
显然不成立。

\paragraph{(b) If P then Q. Not P. Therefore, not
Q.}\label{b-if-p-then-q.-not-p.-therefore-not-q.}

这是\textbf{否前式} (denying the antecedent)。\\
反例：\\
如果我在北京，则我在中国。\\
我不在北京。\\
所以我不在中国。\\
显然错误。

\subsubsection{Exercise
8：证明等价}\label{exercise-8ux8bc1ux660eux7b49ux4ef7}

\begin{enumerate}
\def\labelenumi{\arabic{enumi}.}
\tightlist
\item
  P → Q 与 ¬P ∨ Q：等价。\\
\item
  ¬(P ∨ Q) 与 ¬P ∧ ¬Q：等价。\\
\item
  ¬(P ∧ Q) 与 ¬P ∨ ¬Q：等价。
\end{enumerate}

这些都属于期中必须熟练掌握的标准公式。

\subsubsection{Exercise
9：否定并化简}\label{exercise-9ux5426ux5b9aux5e76ux5316ux7b80}

\begin{enumerate}
\def\labelenumi{\arabic{enumi}.}
\tightlist
\item
  ¬(P ∧ Q) ≡ ¬P ∨ ¬Q
\item
  ¬(P ∨ Q) ≡ ¬P ∧ ¬Q
\item
  ¬(P → Q) ≡ P ∧ ¬Q
\item
  ¬(P ↔ Q) ≡ (P ∧ ¬Q) ∨ (¬P ∧ Q)
\end{enumerate}

\begin{center}\rule{0.5\linewidth}{0.5pt}\end{center}

\subsection{1.9 Homework 2 详解（Predicate Logic and
Quantifiers）}\label{homework-2-ux8be6ux89e3predicate-logic-and-quantifiers}

设论域是``大学中的所有学生''。

\subsubsection{第 1
题：翻译成符号}\label{ux7b2c-1-ux9898ux7ffbux8bd1ux6210ux7b26ux53f7}

设\\
S(x)：x 学习数学\\
P(x)：x 通过考试

\begin{enumerate}
\def\labelenumi{\arabic{enumi}.}
\item
  Every mathematics student passed the exam.\\
  ∀x (S(x) → P(x))
\item
  Some student passed the exam.\\
  ∃x P(x)
\item
  There is a mathematics student who did not pass the exam.\\
  ∃x (S(x) ∧ ¬P(x))
\end{enumerate}

\subsubsection{第 2
题：翻译成英文}\label{ux7b2c-2-ux9898ux7ffbux8bd1ux6210ux82f1ux6587}

\begin{enumerate}
\def\labelenumi{\arabic{enumi}.}
\item
  ∀x∃y (x ≠ y)\\
  直译为：对每个 x，都存在某个 y，使得 x ≠ y。\\
  若论域至少有两个元素，这句话为真；其含义是``没有哪个元素是整个论域里唯一的元素''。
\item
  ∃y∀x (x ≤ y)\\
  存在一个 y，使得每个 x 都满足 x ≤ y。\\
  也就是``存在一个最大元素''。
\end{enumerate}

\subsubsection{第 3
题：否定量词句}\label{ux7b2c-3-ux9898ux5426ux5b9aux91cfux8bcdux53e5}

\begin{enumerate}
\def\labelenumi{\arabic{enumi}.}
\tightlist
\item
  ¬∀xP(x) ≡ ∃x¬P(x)
\item
  ¬∃x(P(x) ∧ Q(x)) ≡ ∀x ¬(P(x) ∧ Q(x)) ≡ ∀x(¬P(x) ∨ ¬Q(x))
\end{enumerate}

\subsubsection{第 4
题：判断论证是否有效}\label{ux7b2c-4-ux9898ux5224ux65adux8bbaux8bc1ux662fux5426ux6709ux6548}

前提： - ∀x(P(x) → Q(x)) - ∃xP(x)

结论： - ∃xQ(x)

这是\textbf{有效}的。\\
理由：由第二个前提，取某个 a，使 P(a) 成立；再由第一个前提得 P(a) →
Q(a)，故 Q(a) 成立；于是存在一个 x 使 Q(x) 成立，即 ∃xQ(x)。

\begin{center}\rule{0.5\linewidth}{0.5pt}\end{center}

\subsection{1.10 Worksheet 2 详解（Predicate Logic in
Practice）}\label{worksheet-2-ux8be6ux89e3predicate-logic-in-practice}

\subsubsection{Exercise 1：Prime Numbers}\label{exercise-1prime-numbers}

题目给出\\
``n 是素数''： ∀d ∈ N ( d \textbar{} n ⇒ (d = 1 ∨ d = n) )

\paragraph{(a) 写出否定}\label{a-ux5199ux51faux5426ux5b9a}

¬∀d∈N (d\textbar n ⇒ (d=1 ∨ d=n)) ≡ ∃d∈N ¬(d\textbar n ⇒ (d=1 ∨ d=n)) ≡
∃d∈N (d\textbar n ∧ ¬(d=1 ∨ d=n)) ≡ ∃d∈N (d\textbar n ∧ d≠1 ∧ d≠n)

\paragraph{(b) 英文解释}\label{b-ux82f1ux6587ux89e3ux91ca}

``存在一个自然数 d，它整除 n，但 d 既不是 1 也不是 n。''

\paragraph{(c) 为什么这就是 not
prime}\label{c-ux4e3aux4ec0ux4e48ux8fd9ux5c31ux662f-not-prime}

因为素数的定义正是``除了 1 和自身外没有别的因子''；\\
它的否定自然就是``有一个非平凡因子''。

\subsubsection{Exercise 2：Course
Prerequisites}\label{exercise-2course-prerequisites}

原命题： ∀x ((Student(x) ∧ Enrolls(x)) ⇒ Prereq(x))

\paragraph{(a) 否定}\label{a-ux5426ux5b9a}

∃x ¬((Student(x) ∧ Enrolls(x)) ⇒ Prereq(x)) ≡ ∃x ((Student(x) ∧
Enrolls(x)) ∧ ¬Prereq(x))

\paragraph{(b) 英文解释}\label{b-ux82f1ux6587ux89e3ux91ca-1}

``存在某个学生，他选修了这门课，但并未满足先修条件。''

\paragraph{(c)
一个具体例子}\label{c-ux4e00ux4e2aux5177ux4f53ux4f8bux5b50}

例如某学生没有修过 Calculus I，却注册了 Calculus II。

\subsubsection{Exercise 3：Library Books}\label{exercise-3library-books}

\begin{enumerate}
\def\labelenumi{(\Alph{enumi})}
\tightlist
\item
  ∀b (Book(b) ⇒ ∃x Borrows(x,b))\\
\item
  ∃x ∀b (Book(b) ⇒ Borrows(x,b))
\end{enumerate}

\paragraph{(a) 含义}\label{a-ux542bux4e49}

\begin{itemize}
\tightlist
\item
  (A)：每一本书都至少被某个人借过。
\item
  (B)：存在一个人，把所有书都借了。
\end{itemize}

\paragraph{(b) 哪个更合理}\label{b-ux54eaux4e2aux66f4ux5408ux7406}

现实中 (A) 更可能合理；(B) 往往太强了。

\paragraph{(c) 一个使 (A) 真而 (B)
假的情景}\label{c-ux4e00ux4e2aux4f7f-a-ux771fux800c-b-ux5047ux7684ux60c5ux666f}

图书馆有 100 本书，每本书都被不同学生借走过；但没有任何一个学生借过全部
100 本。\\
则 (A) 真，(B) 假。

\subsubsection{Exercise 4：Diagnosing
Errors}\label{exercise-4diagnosing-errors}

\paragraph{(a) ∃d∈N (d\textbar n ⇒ (d=1 ∨ d=n))}\label{a-dn-dn-d1-dn}

实际意思：\\
``存在一个 d，使得如果 d 整除 n，那么 d 是 1 或 n。''

为什么错：\\
取 d=2，如果 2 根本不整除
n，则蕴含自动为真，于是该式几乎总是真，完全没有表达素数定义。

正确写法：\\
∀d∈N (d\textbar n ⇒ (d=1 ∨ d=n))

\paragraph{(b) ∀x∃y Attends(x,y)}\label{b-xy-attendsxy}

若 intended 是``Every lecture is attended by at least one
student''，则此式很可能把 x 当学生、y 当 lecture，结果变成\\
``每个人都参加某个活动/讲座''，顺序不对。

正确写法之一：\\
∀y (Lecture(y) → ∃x (Student(x) ∧ Attends(x,y)))

\paragraph{(c) ∀a∃d SubmittedBefore(a,d)}\label{c-ad-submittedbeforead}

实际意思：\\
``每个作业 a 都有某个（可能不同的）截止日期 d，使得 a 在 d 之前提交。''

为什么错：\\
原意要表达``存在同一个统一 deadline''，而这里 d 可以随 a 改变。

正确写法：\\
∃d ∀a SubmittedBefore(a,d)

\subsubsection{Exercise 5：English to
Logic}\label{exercise-5english-to-logic}

下面只给一种合理写法，实际可有多种。

\begin{enumerate}
\def\labelenumi{\arabic{enumi}.}
\item
  Every student has taken at least one mathematics course.\\
  ∀x (Student(x) → ∃c (MathCourse(c) ∧ Took(x,c)))
\item
  There is a professor who teaches every course in the department.\\
  ∃p (Professor(p) ∧ ∀c (Course(c) → Teaches(p,c)))
\item
  Every exam has a question that every student finds difficult.\\
  ∀e (Exam(e) → ∃q (QuestionOf(q,e) ∧ ∀s (Student(s) →
  DifficultFor(s,q))))
\end{enumerate}

\subsubsection{Exercise 6：Logic to
English}\label{exercise-6logic-to-english}

\begin{enumerate}
\def\labelenumi{\arabic{enumi}.}
\item
  ∀x∃y (x ≠ y ∧ Friend(x,y))\\
  每个人都有一个与自己不同的朋友。
\item
  ∃x∀y (Student(y) ⇒ Older(x,y))\\
  存在某个人，他比所有学生都年长。
\end{enumerate}

\begin{center}\rule{0.5\linewidth}{0.5pt}\end{center}

\section{第二章 集合、函数与关系（Sets, Functions,
Relations）}\label{ux7b2cux4e8cux7ae0-ux96c6ux5408ux51fdux6570ux4e0eux5173ux7cfbsets-functions-relations}

\subsection{2.1 核心总结}\label{ux6838ux5fc3ux603bux7ed3-1}

这一章是整个课程最``语言基础''的部分。后面从 N 构造 Z、从 Z 构造
Q，全都要靠这里的概念完成。\\
你必须对下面五类对象特别熟悉：

\begin{enumerate}
\def\labelenumi{\arabic{enumi}.}
\tightlist
\item
  集合与子集；
\item
  集合运算；
\item
  函数；
\item
  关系；
\item
  偏序与等价关系。
\end{enumerate}

\subsubsection{本章必背}\label{ux672cux7ae0ux5fc5ux80cc}

\begin{itemize}
\tightlist
\item
  A = B 当且仅当 ∀x (x∈A ↔ x∈B)。
\item
  A ⊂ B 且 B ⊂ A ⇒ A = B。
\item
  A ∩ (B ∪ C) = (A ∩ B) ∪ (A ∩ C)。
\item
  (A ∪ B)\^{}c = A\^{}c ∩ B\^{}c，(A ∩ B)\^{}c = A\^{}c ∪ B\^{}c。
\item
  函数是 X×Y 的特殊子集：对每个 x∈X，存在唯一 y∈Y 使 (x,y)∈f。
\item
  f 可逆 ⇔ f 是双射 (bijection)。
\item
  偏序：自反 + 反对称 + 传递。
\item
  等价关系：自反 + 对称 + 传递。
\item
  等价类把集合分成互不相交的若干块；商集是``所有等价类''的集合。
\end{itemize}

\subsection{2.2 集合（Sets）与集合代数（Set
Algebra）概念整理}\label{ux96c6ux5408setsux4e0eux96c6ux5408ux4ee3ux6570set-algebraux6982ux5ff5ux6574ux7406}

\subsubsection{Axiom of
Extension（外延公理）}\label{axiom-of-extensionux5916ux5ef6ux516cux7406}

两个集合相等，只看元素，不看写法： A = B ⇔ ∀x (x∈A ↔ x∈B)

所以证明集合相等，最标准的方法是： - 先证 A ⊂ B； - 再证 B ⊂ A。

\subsubsection{常见集合构造}\label{ux5e38ux89c1ux96c6ux5408ux6784ux9020}

\begin{itemize}
\tightlist
\item
  并集：A ∪ B = \{x \textbar{} x∈A 或 x∈B\}
\item
  交集：A ∩ B = \{x \textbar{} x∈A 且 x∈B\}
\item
  差集：\(A\setminus B=\{x\mid x\in A   ext{ 且 } x
  otin B\}\)
\item
  笛卡尔积：A × B = \{(a,b) \textbar{} a∈A, b∈B\}
\item
  幂集：P(A) = \{X \textbar{} X ⊂ A\}
\end{itemize}

\subsubsection{集合恒等式的常见证明方式}\label{ux96c6ux5408ux6052ux7b49ux5f0fux7684ux5e38ux89c1ux8bc1ux660eux65b9ux5f0f}

证明 A ∩ (B ∪ C) = (A ∩ B) ∪ (A ∩ C) 时，一般写：

取任意 x。\\
x ∈ A ∩ (B ∪ C)\\
⇔ x∈A 且 (x∈B 或 x∈C)\\
⇔ (x∈A 且 x∈B) 或 (x∈A 且 x∈C)\\
⇔ x ∈ (A ∩ B) ∪ (A ∩ C)

这就是``元素追踪法 (element chasing)''。

\subsection{2.3
函数（Functions）概念整理}\label{ux51fdux6570functionsux6982ux5ff5ux6574ux7406}

\subsubsection{set-theoretic definition}\label{set-theoretic-definition}

函数 f : X → Y 在集合论中被定义为 X×Y 的一个子集 f，满足：

对每个 x∈X，存在唯一 y∈Y，使 (x,y)∈f。

也就是说： - 不能漏掉任何 x； - 每个 x 不能对应多个输出。

\subsubsection{image / preimage}\label{image-preimage}

\begin{itemize}
\tightlist
\item
  f(x)：元素 x 的像 (image)
\item
  f(A)：子集 A 的像
\item
  f(X)：整个函数的像，也就是实际输出集合
\item
  \(f^{-1}(B)\)：\(B\) 的原像 (preimage)
\end{itemize}

非常重要：\\
\textbf{原像永远可定义}，即使 f 没有逆函数。\\
这是很多同学最容易混淆的地方。

\subsubsection{injective / surjective /
bijective}\label{injective-surjective-bijective}

\begin{itemize}
\tightlist
\item
  单射 (injective)：不同输入不会撞车；
\item
  满射 (surjective)：目标集合每个元素都能被打到；
\item
  双射 (bijective)：既单又满，可逆。
\end{itemize}

\subsubsection{inverse function}\label{inverse-function}

逆函数存在当且仅当： - 有左逆 (left inverse)； - 有右逆 (right
inverse)； - 等价地，f 是双射。

\begin{center}\rule{0.5\linewidth}{0.5pt}\end{center}

\subsection{2.4 Lecture Notes 习题详解（Exercises
11--24）}\label{lecture-notes-ux4e60ux9898ux8be6ux89e3exercises-1124}

\subsubsection{Exercise 11}\label{exercise-11}

\textbf{题意}：证明 A ⊂ B 且 B ⊂ A 蕴含 A = B。

\textbf{解答}：\\
由 A ⊂ B，得任意 x∈A 都有 x∈B；\\
由 B ⊂ A，得任意 x∈B 都有 x∈A。\\
因此对任意 x，x∈A ↔ x∈B。\\
由外延公理 (Axiom of Extension)，A = B。

\subsubsection{Exercise 12}\label{exercise-12}

\textbf{题意}：若 \textbar A\textbar=5，\textbar B\textbar=3，求
\textbar A×B\textbar。

\textbf{解答}：\\
A 中每个元素都可以与 B 中 3 个元素配对；\\
总数为 5×3=15。\\
所以 \textbar A×B\textbar{} = 15。

\subsubsection{Exercise 13}\label{exercise-13}

\textbf{题意}：写出 intersection 的类似性质并证明。

\textbf{解答}：与 Lemma 1 对应，可写为：

\begin{itemize}
\tightlist
\item
  A ∩ E = A（若 E 是全集）
\item
  A ∩ B = B ∩ A（交换律）
\item
  A ∩ (B ∩ C) = (A ∩ B) ∩ C（结合律）
\item
  A ∩ A = A（幂等律）
\item
  A ⊂ B 当且仅当 A ∩ B = A
\end{itemize}

最后一条证明：\\
若 A⊂B，则 A 的每个元素都在 B 中，所以 A∩B 就正好是 A。\\
反过来若 A∩B=A，则任意 x∈A 都在 A∩B 中，所以 x∈B，故 A⊂B。

\subsubsection{Exercise 14}\label{exercise-14}

\textbf{题意}：判断\\
(A ∩ B) ∪ C = A ∩ (B ∪ C)\\
是否当且仅当 C ⊂ A。

\textbf{解答}：这个命题为真。

先看右边： A ∩ (B ∪ C) = (A ∩ B) ∪ (A ∩ C)

因此原式等价于： (A ∩ B) ∪ C = (A ∩ B) ∪ (A ∩ C)

若 C ⊂ A，则 A∩C = C，所以上式成立。\\
反过来若上式成立，取任意 x∈C，则左边 x 在 (A∩B)∪C 中，因此也在右边。若 x
不在 A∩B，则必须在 A∩C 中，因而 x∈A。\\
所以 C⊂A。\\
故``当且仅当''成立。

\subsubsection{Exercise 15}\label{exercise-15}

\textbf{题意}：证明所有从 A 到 B 的函数组成一个集合 B\^{}A。

\textbf{解答}：\\
A×B 是集合，所以 P(A×B) 也是集合。\\
每个函数 f : A→B 都可看作 A×B 的某个子集。\\
于是我们可用分离公理 / comprehension 从 P(A×B) 中选出那些满足``对每个
a∈A，存在唯一 b∈B 使 (a,b)∈f''的子集。\\
因此所有这类 f 的全体构成一个集合，记为 B\^{}A。

\subsubsection{Exercise 16}\label{exercise-16}

\textbf{题意}：举一些函数例子。

\textbf{解答}：

\begin{itemize}
\tightlist
\item
  f : Z → Z, f(n)=n+1
\item
  g : R → R, g(x)=x\^{}2
\item
  h : \{a,b,c\} → \{0,1\}, 设 h(a)=0, h(b)=1, h(c)=1
\item
  \(id_X : X\to X,\ x\mapsto x\)
\end{itemize}

\subsubsection{Exercise 17}\label{exercise-17}

\textbf{题意}：举单射/满射的例子与反例。

\textbf{解答}：

\begin{enumerate}
\def\labelenumi{\arabic{enumi}.}
\item
  f : Z→Z, f(n)=n+1\\
  单射且满射，因此双射。
\item
  g : R→R, g(x)=x\^{}2\\
  不是单射（1 与 -1 撞车），也不是满射（负数打不到）。
\item
  h : R→{[}0,∞), h(x)=x\^{}2\\
  不是单射，但满射。
\item
  k : R→R, k(x)=e\^{}x\\
  单射，但不是满射（≤0 的数打不到）。
\end{enumerate}

\subsubsection{Exercise 18}\label{exercise-18}

\textbf{题意}：证明逆函数若存在则唯一。

\textbf{解答}：\\
设 g,h : Y→X 都是 f 的逆，即 \(g\circ f = id_X,\ f\circ g = id_Y\)\\
\(h\circ f = id_X,\ f\circ h = id_Y\)

则 \(g = g\circ id_Y\) = g\circ (f\circ h) = (g\circ f)\circ h =
id\_X\circ h = h

所以逆函数唯一。

\subsubsection{Exercise 19}\label{exercise-19}

\textbf{题意}：证明 f 有逆函数当且仅当 f 是双射。

\textbf{解答}：

\paragraph{必要性}\label{ux5fc5ux8981ux6027}

若 f 有逆 g，则 - 若 \(f(x_1)=f(x_2)\)，两边作用 \(g\)，得
\(x_1=x_2\)，所以 \(f\) 单射； - 对任意 y∈Y，y=f(g(y))，所以 y 在 f
的像中，因此 f 满射。

故 f 为双射。

\paragraph{充分性}\label{ux5145ux5206ux6027}

若 f 为双射，则对每个 y∈Y： - 因为满射，存在 x 使 f(x)=y； -
因为单射，这个 x 唯一。

于是定义 g(y)=唯一满足 f(x)=y 的 x。\\
即可得 \(g=f^{-1}\)，且满足左右逆条件。

\subsubsection{Exercise 20}\label{exercise-20}

\textbf{题意}：对关系 y\^{}2 = x（定义在 Z×Z 上）求 domain 和 image。

\textbf{解答}：

\begin{itemize}
\tightlist
\item
  domain：所有能写成某个整数平方的整数，即\\
  \{0,1,4,9,16,\ldots\}
\item
  image：所有整数，因为任给 y∈Z，取 x=y\^{}2，就有 y\^{}2=x。\\
  所以 image = Z。
\end{itemize}

\subsubsection{Exercise 21}\label{exercise-21}

\textbf{题意}：举一些满足 / 不满足各种关系性质的例子。

\textbf{解答}：

\begin{itemize}
\item
  自反 (reflexive) 例子：x≤y 定义在 Z 上中的``≤''是自反的，因为 x≤x。
\item
  非自反例子：x\textless y 不是自反的，因为 x\textless x 永假。
\item
  对称 (symmetric) 例子：xRy 若 \textbar x-y\textbar≤1，是对称的。
\item
  非对称例子：x≤y 不是对称的。
\item
  反对称 (antisymmetric) 例子：x≤y 是反对称的。
\item
  非反对称例子：在实数上定义 xRy 若 \textbar x-y\textbar≤1，则 0R1 且
  1R0 但 0≠1。
\item
  传递 (transitive) 例子：≤、整除关系 \textbar{} 都传递。
\item
  非传递例子：xRy 若 \textbar x-y\textbar≤1，则 0R1、1R2，但 0R2
  不成立。
\end{itemize}

\subsubsection{Exercise 22}\label{exercise-22}

\textbf{题意}：说明给定关系不是 partial order，失败的是哪条公理。

\textbf{解答}：\\
由于定义中额外规定了 \(1\preceq 3\) 且 \(3\preceq 1\)，但
\(1\neq 3\)，\\
故\textbf{反对称性}失败。\\
所以它不是偏序 (partial order)。

\subsubsection{Exercise 23}\label{exercise-23}

\textbf{题意}：分别举出``只坏掉一个性质，另外两个成立''的例子。

\textbf{解答}：

\begin{enumerate}
\def\labelenumi{\arabic{enumi}.}
\item
  \textbf{不自反，但对称且传递}\\
  在集合 \{1,2\} 上定义 R = \{(1,1)\}。\\
  对元素 2 不满足 2R2，所以不自反；\\
  同时它对称且传递。
\item
  \textbf{自反且传递，但不对称}\\
  取实数上的 ≤。\\
  自反、传递都成立，但一般 x≤y 不推出 y≤x。
\item
  \textbf{自反且对称，但不传递}\\
  在整数上定义 xRy 当且仅当 \textbar x-y\textbar≤1。\\
  则 xRx，且对称；\\
  但 0R1 与 1R2 成立，而 0R2 不成立。
\end{enumerate}

\subsubsection{Exercise 24}\label{exercise-24}

\textbf{题意}：证明模 m 同余是等价关系。

\textbf{解答}：定义\\
a ≡ b (mod m) ⇔ m \textbar{} (a-b)

要证自反、对称、传递。

\begin{enumerate}
\def\labelenumi{\arabic{enumi}.}
\item
  \textbf{自反}\\
  a-a=0，而 m\textbar0，所以 a≡a (mod m)。
\item
  \textbf{对称}\\
  若 m\textbar(a-b)，则 a-b = km。\\
  于是 b-a = -km，也被 m 整除，所以 b≡a (mod m)。
\item
  \textbf{传递}\\
  若 m\textbar(a-b) 且 m\textbar(b-c)，则\\
  a-b = km, b-c = ℓ m。\\
  相加得 a-c = (k+ℓ)m，故 m\textbar(a-c)，所以 a≡c (mod m)。
\end{enumerate}

因此它是等价关系。

\begin{center}\rule{0.5\linewidth}{0.5pt}\end{center}

\subsection{2.5 Worksheet 3 详解（Sets and
Functions）}\label{worksheet-3-ux8be6ux89e3sets-and-functions}

\subsubsection{Exercise 1：Hiking trip}\label{exercise-1hiking-trip}

总人数 10。\\
7 人涂防晒，6 人戴帽子，2 人完全没有防护。\\
所以有防护的人数为 10-2=8。\\
由容斥原理： \textbar S∩H\textbar{} = \textbar S\textbar{} +
\textbar H\textbar{} - \textbar S∪H\textbar{} = 7 + 6 - 8 = 5

答案：\textbf{5 人}同时涂防晒且戴帽子。

\subsubsection{Exercise 2：Venn diagram}\label{exercise-2venn-diagram}

这题主要考你能否把条件翻译成图形关系。

\begin{enumerate}
\def\labelenumi{\arabic{enumi}.}
\item
  A⊂B，C⊂B，A∩C=∅\\
  画法：A 和 C 都完全放在 B 里面，且 A、C 分开不重叠。
\item
  A∩B≠∅，A∩C≠∅，B∩C≠∅，A∩B∩C=∅\\
  画法：三个集合两两有交，但中间三重交区域留空。
\item
  A⊂(B∩C)，B⊂C\\
  画法：C 最大，B 完全在 C 内，A 又完全在 B 与 C 的重合部分内；由于
  B⊂C，所以其实 A⊂B⊂C。
\item
  A⊂(B∪C)，但 A 既不是 B 的子集，也不是 C 的子集\\
  画法：A 横跨 B 与 C，A 的一部分在 \(B\setminus C\)，另一部分在
  \(C\setminus B\)，也可能有一部分在 \(B\cap C\)；关键是 A
  不能完全落在某一个集合中。
\end{enumerate}

\subsubsection{Exercise 3：Left and right
inverses}\label{exercise-3left-and-right-inverses}

设 X=\{a,b,c\}，Y=\{α,β,γ,δ\}。

\paragraph{\texorpdfstring{(1) 构造
\(f_1 : X\to Y\)，有左逆但无右逆}{(1) 构造 f\_1 : X\textbackslash to Y，有左逆但无右逆}}\label{ux6784ux9020-f_1-xto-yux6709ux5de6ux9006ux4f46ux65e0ux53f3ux9006}

取 \(f_1(a)=\alpha,\ f_1(b)=\beta,\ f_1(c)=\gamma\)

它是单射但不是满射（δ 没被打到）。\\
因此有左逆，没有右逆。\\
可定义左逆 h : Y→X 为 h(α)=a, h(β)=b, h(γ)=c，h(δ) 随便定义成 a。\\
则 \(h\circ f_1 = id_X\)。

\paragraph{\texorpdfstring{(2) 构造
\(f_2 : Y\to X\)，有右逆但无左逆}{(2) 构造 f\_2 : Y\textbackslash to X，有右逆但无左逆}}\label{ux6784ux9020-f_2-yto-xux6709ux53f3ux9006ux4f46ux65e0ux5de6ux9006}

取 \(f_2(\alpha)=a,\ f_2(\beta)=a,\ f_2(\gamma)=b,\ f_2(\delta)=c\)

这是满射但不是单射。\\
因此有右逆没有左逆。\\
可定义 g : X→Y 为 g(a)=α, g(b)=γ, g(c)=δ。\\
则 \(f_2\circ g = id_X\)。

\paragraph{(3) 证明若 g : X→X 有左逆 h，则 h
也是右逆}\label{ux8bc1ux660eux82e5-g-xx-ux6709ux5de6ux9006-hux5219-h-ux4e5fux662fux53f3ux9006}

已知 \(h\circ g = id_X\)。\\
这说明 g 是单射。\\
因为 g : X→X 是从有限集合到自身的单射，所以也是满射。\\
于是 \(g\) 可逆，而逆函数唯一；左逆 \(h\) 必等于真正的逆，因此也满足
\(g\circ h = id_X\)。\\
所以 h 同时也是右逆。

\subsubsection{Exercise 4：N 与 Z
的双射}\label{exercise-4n-ux4e0e-z-ux7684ux53ccux5c04}

构造 f : N→Z：

\begin{itemize}
\tightlist
\item
  f(0)=0
\item
  对 n≥1：

  \begin{itemize}
  \tightlist
  \item
    若 n=2k-1 为奇数，则 f(n)=k
  \item
    若 n=2k 为偶数，则 f(n)=-k
  \end{itemize}
\end{itemize}

写成前几项： \(0 \mapsto 0\)\\
\(1 \mapsto 1\)\\
\(2 \mapsto -1\)\\
\(3 \mapsto 2\)\\
\(4 \mapsto -2\)\\
\(5 \mapsto 3\)\\
\(6 \mapsto -3,\ \dots\)

这显然覆盖全部整数，且每个整数恰好出现一次，因此是双射。

\subsubsection{Exercise 5（Honours）：构造注入 f :
N×N→N}\label{exercise-5honoursux6784ux9020ux6ce8ux5165-f-nnn}

经典做法之一是 Cantor pairing function： f(m,n) = ((m+n)(m+n+1))/2 + n

它把每个有序对映到一个唯一自然数，因此是单射。\\
如果考试不要求 honours 细节，你至少要知道：\\
``通过按对角线枚举，也可以把 N×N 注入甚至双射到 N 中''。

\section{第三章 自然数、整数、有理数（Naturals, Integers,
Rationals）}\label{ux7b2cux4e09ux7ae0-ux81eaux7136ux6570ux6574ux6570ux6709ux7406ux6570naturals-integers-rationals}

\subsection{3.1 核心总结}\label{ux6838ux5fc3ux603bux7ed3-2}

这一章最重要的思想不是``数字怎么计算''，而是：

\begin{itemize}
\tightlist
\item
  我们可以从更原始的结构\textbf{构造}数系；
\item
  每一步构造都依赖某种\textbf{等价关系 (equivalence relation)}；
\item
  算术运算必须先证明\textbf{良定义 (well-defined)}。
\end{itemize}

\subsubsection{本章必背}\label{ux672cux7ae0ux5fc5ux80cc-1}

\begin{enumerate}
\def\labelenumi{\arabic{enumi}.}
\tightlist
\item
  自然数模型由 Peano axioms（皮亚诺公理）刻画。\\
\item
  归纳法 (principle of induction) 是自然数最核心的性质。\\
\item
  整数 Z 由 N² 上的等价关系构造： (a,b) \textasciitilde\_Z (c,d) ⇔ a+d =
  b+c
\item
  有理数 Q 由 \(Z\times (Z\setminus\{0\})\) 上的等价关系构造： (a,b)
  \textasciitilde\_Q (c,d) ⇔ ad = bc
\item
  在商集上定义运算时，一定要证明``不依赖代表元''，即 well-defined。
\item
  √2 不是有理数，说明 Q 不足以容纳所有``自然''方程的解。
\end{enumerate}

\subsection{3.2 自然数（Natural Numbers）与 Peano
公理}\label{ux81eaux7136ux6570natural-numbersux4e0e-peano-ux516cux7406}

\subsubsection{Peano
axioms（皮亚诺公理）}\label{peano-axiomsux76aeux4e9aux8bfaux516cux7406}

一个自然数模型由三样东西组成： - 一个集合 N - 一个特殊元素 0∈N -
一个后继函数 S : N→N

满足：

\begin{enumerate}
\def\labelenumi{\arabic{enumi}.}
\tightlist
\item
  \textbf{S 单射}：S(x)=S(y) ⇒ x=y\\
\item
  \textbf{没有不动点}：S(x)≠x\\
\item
  \textbf{0 不是任何元素的后继}：不存在 x 使 S(x)=0\\
\item
  \textbf{归纳原理}：若 P(0) 真，且 P(x)⇒P(S(x))，则 ∀x P(x)
\end{enumerate}

\subsubsection{加法与乘法的递归定义}\label{ux52a0ux6cd5ux4e0eux4e58ux6cd5ux7684ux9012ux5f52ux5b9aux4e49}

加法： - a+0 = a - a+S(b) = S(a+b)

乘法： - a·0 = 0 - a·S(b) = a + (a·b)

这两条定义的精神都是：\\
\textbf{把复杂对象降到更小的对象，再递归返回。}

\subsection{3.3
归纳法（Induction）证明模板}\label{ux5f52ux7eb3ux6cd5inductionux8bc1ux660eux6a21ux677f}

要证明``对所有 n∈N，P(n) 成立''，标准写法：

\begin{enumerate}
\def\labelenumi{\arabic{enumi}.}
\tightlist
\item
  \textbf{Base case}：证明 P(0)
\item
  \textbf{Induction hypothesis}：假设 P(n) 成立
\item
  \textbf{Induction step}：在该假设下证明 P(S(n))
\item
  由归纳原理得 ∀n P(n)
\end{enumerate}

\subsubsection{常见误区}\label{ux5e38ux89c1ux8befux533a}

\begin{itemize}
\tightlist
\item
  只验证几个例子不等于证明；
\item
  归纳步骤不是``从 n 推到 n+1 看起来像对''，而是必须写出明确公式推导；
\item
  若命题中有两个变量，常常是\textbf{固定一个，对另一个做归纳}。
\end{itemize}

\begin{center}\rule{0.5\linewidth}{0.5pt}\end{center}

\subsection{3.4 Lecture Notes 习题详解（Exercises
25--38）}\label{lecture-notes-ux4e60ux9898ux8be6ux89e3exercises-2538}

\subsubsection{Exercise 25}\label{exercise-25}

\textbf{题意}：集合 \{0,1,2\} 配上 S(0)=1, S(1)=2,
S(2)=0，为何不满足全部 Peano 公理？

\textbf{解答}：逐条检查：

\begin{enumerate}
\def\labelenumi{\arabic{enumi}.}
\tightlist
\item
  \textbf{单射}：成立。因为 1、2、0 三个像互不相同。
\item
  \textbf{没有不动点}：成立。S(0)=1≠0，S(1)=2≠1，S(2)=0≠2。
\item
  \textbf{0 不是任何元素的后继}：失败。因为 S(2)=0。
\item
  \textbf{归纳原理}：实际上在这个三元循环里，从 0 出发经 S 可到
  1、2，又回到 0；若 P(0) 成立且
  P(x)⇒P(S(x))，就会推出三者都成立，所以这里归纳性质并不一定坏。
\end{enumerate}

因此至少第三条失败，所以它不是自然数模型。

\subsubsection{Exercise 26}\label{exercise-26}

\textbf{题意}：说明加法的递归定义是``形式上合理''的。

\textbf{解答}：\\
对固定的 a，考虑命题 P(b)：``a+b 已经被唯一地确定''。

\begin{itemize}
\tightlist
\item
  当 b=0 时，由定义 a+0=a，故 P(0) 成立；
\item
  若 P(b) 成立，即 a+b 已被唯一确定，则 a+S(b)=S(a+b) 也被唯一确定，所以
  P(S(b)) 成立。
\end{itemize}

由归纳法，P(b) 对所有 b 成立。\\
因此递归公式确实给出了对每个 b 唯一的值，这就是``定义良好''。

\subsubsection{Exercise 27}\label{exercise-27}

\textbf{题意}：证明对所有 n∈N，0+n=n。

\textbf{解答}：对 n 归纳。

\begin{itemize}
\item
  Base case：n=0\\
  0+0 = 0，成立。
\item
  Induction hypothesis：假设 0+n = n。
\item
  Induction step：\\
  0+S(n) = S(0+n) = S(n)
\end{itemize}

故对所有 n，0+n=n。

\subsubsection{Exercise 28}\label{exercise-28}

\textbf{题意}：计算 2·3。

\textbf{解答}：\\
记 2=S(S(0))，3=S(S(S(0)))。\\
由乘法递归定义：

2·3 = 2·S(2) = 2 + 2·2\\
而 2·2 = 2·S(1)=2+2·1\\
2·1 = 2·S(0)=2+2·0 = 2+0 = 2

所以 2·2 = 2+2 = 4\\
2·3 = 2+4 = 6

\subsubsection{Exercise 29}\label{exercise-29}

\textbf{题意}：从定义推出加法和乘法的交换律、结合律。

\textbf{解答}：这题严格写会很长，考试中你至少要会写出``归纳对象''和关键一步。

\paragraph{（一）加法交换律}\label{ux4e00ux52a0ux6cd5ux4ea4ux6362ux5f8b}

要证：∀a,b，a+b=b+a。\\
对 b 归纳。

\begin{itemize}
\item
  Base case：b=0\\
  a+0=a，而 0+a=a（Exercise 27），故 a+0=0+a。
\item
  Induction hypothesis：假设对某个 b 有 ∀a，a+b=b+a。
\item
  Induction step：\\
  a+S(b)=S(a+b)=S(b+a)=S(b)+a\\
  其中最后一步用到 Worksheet 4 Exercise 1 的结论：S(b)+a = S(b+a)。
\end{itemize}

故交换律成立。

\paragraph{（二）加法结合律}\label{ux4e8cux52a0ux6cd5ux7ed3ux5408ux5f8b}

要证：a+(b+c)=(a+b)+c。\\
对 c 归纳。

\begin{itemize}
\tightlist
\item
  c=0：\\
  a+(b+0)=a+b=(a+b)+0
\item
  若对 c 成立，则\\
  a+(b+S(c)) = a+S(b+c) = S(a+(b+c)) = S((a+b)+c) = (a+b)+S(c)
\end{itemize}

故成立。

\paragraph{（三）乘法交换律}\label{ux4e09ux4e58ux6cd5ux4ea4ux6362ux5f8b}

可对 b 归纳，利用递归式 a·S(b)=a+(a·b)
以及已经证明的加法交换、结合与分配性质逐步推出。\\
考试时若要求详细证明，写法与加法交换律完全同型，只是代数展开更多。

\paragraph{（四）乘法结合律}\label{ux56dbux4e58ux6cd5ux7ed3ux5408ux5f8b}

证明 a·(b·c) = (a·b)·c 也可对 c 归纳，用乘法递归定义与分配律配合完成。

\textbf{考试建议}：\\
这类题比起``记答案''，更重要的是记住： - 对哪个变量归纳； - base case
怎样写； - induction step 中如何用递归定义把 S(c) 拉出来。

\subsubsection{Exercise 30}\label{exercise-30}

\textbf{题意}：证明分配律\\
a·(b+c)=a·b+a·c。

\textbf{解答}：对 c 归纳。

\begin{itemize}
\item
  Base case：c=0\\
  a·(b+0)=a·b\\
  而 a·b + a·0 = a·b + 0 = a·b
\item
  Induction hypothesis：\\
  假设 a·(b+c)=a·b+a·c
\item
  Induction step：\\
  a·(b+S(c)) = a·S(b+c) = a + a·(b+c) = a + (a·b + a·c) = a·b + (a +
  a·c) = a·b + a·S(c)
\end{itemize}

故成立。

\subsubsection{Exercise 31}\label{exercise-31}

\textbf{题意}：说明集合 M=\{0,1,2,\ldots\}∪\{a,b,c\}
满足除归纳法外其他条件。

\textbf{解答}：

\begin{itemize}
\tightlist
\item
  对数值部分，S(n)=n+1；
\item
  对额外元素，S(a)=b, S(b)=c, S(c)=a。
\end{itemize}

检查：

\begin{enumerate}
\def\labelenumi{\arabic{enumi}.}
\tightlist
\item
  \textbf{单射}：每个元素像不同，成立。
\item
  \textbf{没有不动点}：S(x)≠x，成立。
\item
  \textbf{0 不是任何元素后继}：数值部分显然不行；a,b,c 的像是
  b,c,a，也不是 0，成立。
\item
  \textbf{归纳法失败}：令 P(x) 表示``x 是普通数字 0,1,2,\ldots{}
  之一''。\\
  则 P(0) 真；若 P(n) 真，则 P(S(n)) 也真。\\
  但是 P(a)、P(b)、P(c) 都假。\\
  所以``P(0) 成立且 P(x)⇒P(S(x))''并不推出对所有 x 成立。\\
  因而归纳原理失败。
\end{enumerate}

\subsubsection{Exercise 32}\label{exercise-32}

\textbf{题意}：验证 (a,b) \textasciitilde\_Z (c,d) ⇔ a+d=b+c
是等价关系。

\textbf{解答}：

\begin{enumerate}
\def\labelenumi{\arabic{enumi}.}
\item
  \textbf{自反}\\
  (a,b)\textasciitilde(a,b)，因为 a+b=b+a。
\item
  \textbf{对称}\\
  若 a+d=b+c，则交换左右得 c+b=d+a，所以 (c,d)\textasciitilde(a,b)。
\item
  \textbf{传递}\\
  若 (a,b)\textasciitilde(c,d) 且 (c,d)\textasciitilde(e,f)，则
  a+d=b+c\\
  c+f=d+e\\
  两式相加得 a+d+c+f = b+c+d+e 消去 c,d，得 a+f=b+e 所以
  (a,b)\textasciitilde(e,f)。
\end{enumerate}

故 \textasciitilde\_Z 是等价关系。

\subsubsection{Exercise 33}\label{exercise-33}

\textbf{题意}：在 Z 上定义减法。

\textbf{解答}：\\
若 {[}(a,b){]} 表示``a-b''，{[}(c,d){]} 表示``c-d''，\\
那么 (a-b) - (c-d) = (a+d) - (b+c)

所以可定义 {[}(a,b){]} - {[}(c,d){]} := {[}(a+d, b+c){]}

也可以理解为： x-y = x + (-y)

\subsubsection{Exercise 34}\label{exercise-34}

\textbf{题意}：证明整数乘法定义良好。

\textbf{定义}： {[}(a,b){]}·{[}(c,d){]} := {[}(ac+bd, ad+bc){]}

\textbf{解答}：\\
要证：若 (a,b)\textasciitilde(a',b') 且 (c,d)\textasciitilde(c',d')，则
(ac+bd, ad+bc) \textasciitilde{} (a'c'+b'd', a'd'+b'c')

已知 a+b' = b+a'\\
c+d' = d+c'

展开并整理即可得到： (ac+bd) + (a'd'+b'c') = (ad+bc) + (a'c'+b'd')

因此乘法不依赖代表元，故良定义。

\textbf{说明}：考试中这类题最重要的是：\\
- 先写出``要证明什么''；\\
- 再把已知等价条件代入；\\
- 最后通过加法交换结合把左右整理成同一个式子。

\subsubsection{Exercise 35}\label{exercise-35}

\textbf{题意}：证明整数加法与乘法满足交换律、结合律、分配律。

\textbf{解答}：\\
由于整数运算是从自然数运算继承而来，而定义式： - 加法：{[}(a,b){]} +
{[}(c,d){]} = {[}(a+c,b+d){]} - 乘法：{[}(a,b){]}·{[}(c,d){]} =
{[}(ac+bd,ad+bc){]}

本质上只是把自然数中的对应运算搬到等价类层面。\\
例如加法交换律：

{[}(a,b){]} + {[}(c,d){]} = {[}(a+c,b+d){]} = {[}(c+a,d+b){]} =
{[}(c,d){]} + {[}(a,b){]}

其余结合律与分配律也同样通过展开后应用自然数运算的对应性质即可。\\
这里的``严谨点''在于：你已经在 Exercise 34 一类题中确认了这些运算是
well-defined，所以这样的代表元计算是合法的。

\subsubsection{Exercise 36}\label{exercise-36}

\textbf{题意}：证明 \textasciitilde\_Q 是
Y=\(Z\times (Z\setminus\{0\})\) 上的等价关系。

\textbf{定义}： (a,b) \textasciitilde\_Q (c,d) ⇔ ad = bc

\textbf{解答}：

\begin{enumerate}
\def\labelenumi{\arabic{enumi}.}
\tightlist
\item
  \textbf{自反}：ab=ba，所以 (a,b)\textasciitilde(a,b)。
\item
  \textbf{对称}：若 ad=bc，则 cb=da，故 (c,d)\textasciitilde(a,b)。
\item
  \textbf{传递}：\\
  若 ad=bc 且 cf=de，\\
  则 adf = bcf = bde。\\
  因 d≠0，可在整数环里整理为 af=be，故 (a,b)\textasciitilde(e,f)。
\end{enumerate}

所以 \textasciitilde\_Q 是等价关系。

\subsubsection{Exercise 37}\label{exercise-37}

\textbf{题意}：证明有理数乘法定义良好。

\textbf{定义}： {[}(a,b){]}·{[}(c,d){]} := {[}(ac, bd){]}

\textbf{解答}：\\
若 (a,b)\textasciitilde(a',b') 与 (c,d)\textasciitilde(c',d')，即 ab' =
ba'，cd' = dc'。

则 (ac)(b'd') = (ab')(cd') = (ba')(dc') = (a'c')(bd)

故 (ac,bd) \textasciitilde{} (a'c',b'd')。\\
所以乘法良定义。

\subsubsection{Exercise 38}\label{exercise-38}

\textbf{题意}：证明 Q 上加法乘法满足交换律、结合律与分配律。

\textbf{解答}：\\
和整数部分一样，直接从定义展开即可。

例如加法交换： {[}(a,b){]} + {[}(c,d){]} = {[}(ad+bc, bd){]} =
{[}(cb+da, db){]} = {[}(c,d){]} + {[}(a,b){]}

乘法交换： {[}(a,b){]}·{[}(c,d){]} = {[}(ac,bd){]} = {[}(ca,db){]} =
{[}(c,d){]}·{[}(a,b){]}

结合律与分配律的证明都是``展开两边，利用 Z 中运算性质得到相同结果''。

\begin{center}\rule{0.5\linewidth}{0.5pt}\end{center}

\subsection{3.5 Worksheet 4 详解（Naturals and
Integers）}\label{worksheet-4-ux8be6ux89e3naturals-and-integers}

\subsubsection{Exercise 1}\label{exercise-1-1}

\textbf{题意}：证明 S(a)+b = S(a+b)。

\textbf{解答}：对 b 归纳。

\begin{itemize}
\item
  Base case：b=0\\
  S(a)+0 = S(a)\\
  另一方面 S(a+0)=S(a)\\
  所以成立。
\item
  Induction hypothesis：假设 S(a)+b = S(a+b)
\item
  Induction step： S(a)+S(b) = S(S(a)+b)\\
  = S(S(a+b))\\
  = S(a+S(b))
\end{itemize}

故成立。

\subsubsection{Exercise 2}\label{exercise-2-1}

\textbf{题意}：利用上题证明加法交换律。

\textbf{解答}：对 b 归纳，命题为\\
P(b): ∀a, a+b=b+a。

\begin{itemize}
\item
  Base case：b=0\\
  a+0=a=0+a。
\item
  Induction step：假设 ∀a, a+b=b+a。\\
  则 a+S(b)=S(a+b)=S(b+a)=S(b)+a 其中最后一步正是 Exercise 1。
\end{itemize}

故 ∀a,b，a+b=b+a。

\subsubsection{Exercise 3}\label{exercise-3-1}

\textbf{题意}：给出满足除某条 Peano 公理外其余都满足的例子。

\textbf{解答}：

\paragraph{(i) 除 injectivity
外都满足}\label{i-ux9664-injectivity-ux5916ux90fdux6ee1ux8db3}

可取 N=\{0,1\}，定义\\
S(0)=1, S(1)=1。\\
则 S 非单射，因为 S(0)=S(1)。\\
其余条件可适当检查。

\paragraph{(ii) 除 induction
外都满足}\label{ii-ux9664-induction-ux5916ux90fdux6ee1ux8db3}

Lecture Notes Exercise 31 已经给出标准例子：\\
M=\{0,1,2,\ldots\}∪\{a,b,c\}，并让 a,b,c 构成一个循环。\\
这样 0 不在循环里，因此前几条成立，但归纳法失败。

\begin{center}\rule{0.5\linewidth}{0.5pt}\end{center}

\subsection{3.6 Worksheet 5
详解（Rationals）}\label{worksheet-5-ux8be6ux89e3rationals}

\subsubsection{Exercise 1：Relations on Q 是否
well-defined}\label{exercise-1relations-on-q-ux662fux5426-well-defined}

题目给出三个关系，要判断是否与代表元选择无关。

\paragraph{(1) p/q R m/n 当且仅当 p \textgreater{}
m}\label{pq-r-mn-ux5f53ux4e14ux4ec5ux5f53-p-m}

\textbf{不良定义}。\\
因为同一个有理数可写成不同分数。\\
例如 1/2 = 2/4。\\
若拿 1/2 与 1/3 比，则 p=1,m=1，不满足 p\textgreater m；\\
若拿 2/4 与 1/3 比，则 p=2,m=1，满足 p\textgreater m。\\
结果依赖表示法，所以不 well-defined。

\paragraph{(2) p/q R m/n 当且仅当 pn - qm \textgreater{}
0}\label{pq-r-mn-ux5f53ux4e14ux4ec5ux5f53-pn---qm-0}

这是\textbf{良定义}。\\
注意 pn - qm \textgreater{} 0 等价于 pn \textgreater{} qm 等价于 p/q
\textgreater{} m/n
（当分母取正规表示时更直观；严格做法可通过交叉相乘并注意等价类定义）
它本质上只在比较两个有理数大小，和代表元无关。

\paragraph{(3) p/q R m/n 当且仅当 (pn-mq)nq \textgreater{}
0}\label{pq-r-mn-ux5f53ux4e14ux4ec5ux5f53-pn-mqnq-0}

一般\textbf{不良定义}。\\
因为它把 nq 的符号也乘进去了，而分母改写代表时可能改变符号结构。\\
例如把同一个有理数同时写成 1/2 与 -1/-2，会改变 nq
的正负表达方式，从而影响真假。\\
因此不 well-defined。

\subsubsection{Exercise 2：Euclidean
algorithm}\label{exercise-2euclidean-algorithm}

设 a=bk+c，其中 0≤c\textless b。证明 gcd(a,b)=gcd(b,c)。

\textbf{解答}：\\
记 d 是 a,b 的公因子，则 d\textbar a 且 d\textbar b。\\
由 a-bk=c，得 d\textbar c。\\
所以 d 也是 b,c 的公因子。

反过来，若 d\textbar b 且 d\textbar c，则由 a=bk+c，得 d\textbar a。\\
所以 d 也是 a,b 的公因子。

因此 a,b 与 b,c 的公因子集合相同，最大公因子自然相同：\\
gcd(a,b)=gcd(b,c)。

\subsubsection{Exercise 3：Geometric power
series}\label{exercise-3geometric-power-series}

定义 f(n)=1 + 1/2 + 1/2\^{}2 + \ldots{} + 1/2\^{}n

问 f(N) 在 Q 中是否有 supremum。

\textbf{解答}：\\
注意有限等比和公式： f(n)=2 - 1/2\^{}n

所以所有项都小于 2，而且随着 n 增大越来越接近 2。\\
因此 2 是它的上界。\\
并且任意小于 2 的有理数 s，都可写成 s\textless2。令 2-s\textgreater0。\\
因为 1/2\^{}n 可以任意小，所以可找到 n 使 1/2\^{}n \textless{} 2-s，从而
2 - 1/2\^{}n \textgreater{} s。 也就是 f(n)\textgreater s。\\
所以 s 不可能是上界。\\
因此 supremum 存在于 Q 中，且 sup(f(N)) = 2。

\subsubsection{Exercise 4：Rationals are full of
gaps}\label{exercise-4rationals-are-full-of-gaps}

给定 a\textless b，证明可找一个集合 X⊂\{x∈Q \textbar{}
a\textless x\textless b\} 使 sup(X) 在 Q 中不存在。

\textbf{解答}：\\
因为区间 (a,b) 中一定可以找到某个有理数 r，使得 a \textless{} r
\textless{} b/√2 这样的写法不方便；更直接地，取任意有理数 u,v 满足 a
\textless{} u \textless{} v \textless{} b。 考虑集合 X = \{x∈Q
\textbar{} x\^{}2 \textless{} 2 且 u \textless{} x \textless{} v\}
如果区间选得包含``逼近 √2 的有理数们''，那么 X 在 Q
中有上界，但上确界应该是 √2（或某个与 √2 平移缩放后的无理数），不属于
Q。\\
因此它在 Q 中没有 supremum。

一个更标准的构造是：先取有理数 r 使 a \textless{} r \textless{} b 再令 X
= \{x∈Q \textbar{} a \textless{} x \textless{} b 且 x \textless{}
r+√2/N\} 之类并不简洁。\\
最干净的考试写法通常是：

\begin{itemize}
\tightlist
\item
  先在 (a,b) 中取有理数 u\textless v；
\item
  通过线性变换把集合 S=\{q∈Q:q\^{}2\textless2\} 缩放平移到 (u,v) 内；
\item
  其上确界对应于缩放后的 √2，仍不在 Q 中。
\end{itemize}

本题重点不是具体公式，而是理解：\\
\textbf{任意有理区间里都能嵌入一个``缺少 supremum 的有理子集''}。\\
这正说明 Q 到处都有``空隙 (gaps)''。

\begin{center}\rule{0.5\linewidth}{0.5pt}\end{center}

\subsection{3.7 有理数里最重要的一个证明：√2
不是有理数}\label{ux6709ux7406ux6570ux91ccux6700ux91cdux8981ux7684ux4e00ux4e2aux8bc1ux660e2-ux4e0dux662fux6709ux7406ux6570}

这部分虽然在 notes
中已经有完整证明，但期中很可能会考它的结构，因此这里单独整理。

\subsubsection{命题}\label{ux547dux9898}

不存在有理数 x 使 x²=2。

\subsubsection{辅助引理}\label{ux8f85ux52a9ux5f15ux7406}

若整数 p 的平方 p² 为偶数，则 p 为偶数。

\textbf{证明}：用逆否命题。\\
若 p 是奇数，则 p=2k+1，\\
p²=(2k+1)²=4k²+4k+1=2(2k²+2k)+1，仍是奇数。\\
所以若 p² 偶，则 p 不可能奇，只能偶。

\subsubsection{主证明}\label{ux4e3bux8bc1ux660e}

假设存在有理数 x 满足 x²=2。\\
可写 x=p/q 其中 p,q∈Z，q≠0，且 gcd(p,q)=1。

则 p²/q² = 2 所以 p² = 2q²

因此 p² 为偶数，所以 p 为偶数。写 p=2k。\\
代回得 4k² = 2q² 所以 q² = 2k²

于是 q² 为偶数，从而 q 也为偶数。\\
这说明 p,q 都能被 2 整除，与 gcd(p,q)=1 矛盾。\\
故假设错误，√2 不属于 Q。

\subsubsection{这个证明的骨架要记住}\label{ux8fd9ux4e2aux8bc1ux660eux7684ux9aa8ux67b6ux8981ux8bb0ux4f4f}

\begin{enumerate}
\def\labelenumi{\arabic{enumi}.}
\tightlist
\item
  假设是有理数；
\item
  写成最简分数 p/q；
\item
  推出 p 偶；
\item
  再推出 q 偶；
\item
  与最简矛盾。
\end{enumerate}

\begin{center}\rule{0.5\linewidth}{0.5pt}\end{center}

\section{第四章 实数导论（The Reals）至
§4.7}\label{ux7b2cux56dbux7ae0-ux5b9eux6570ux5bfcux8bbathe-realsux81f3-4.7}

\subsection{4.1 核心总结}\label{ux6838ux5fc3ux603bux7ed3-3}

这一章真正想回答的问题是：\\
\textbf{为什么有理数还不够？什么额外性质把实数和有理数区分开？}

答案是：\textbf{完备性 (completeness)}。

Q 是一个有序域 (ordered field)，但它\textbf{不完备}；\\
R 既是有序域，又完备。

\subsubsection{本章必背}\label{ux672cux7ae0ux5fc5ux80cc-2}

\begin{itemize}
\tightlist
\item
  全序 (total order)：任意两个元素都可比较；
\item
  上界 / 下界；
\item
  上确界 / 下确界；
\item
  完备性：每个有上界的子集都有 supremum，每个有下界的子集都有 infimum；
\item
  Q 不完备，反例核心是\\
  S=\{x∈Q \textbar{} x²\textless2\}
\item
  实数模型：有序域 + 完备性；
\item
  §4.7 只是``构造动机''：用越来越精细的近似来刻画一个实数；Dedekind cuts
  本身不考。
\end{itemize}

\subsection{4.2 总序与有序域（Total Orders and Ordered
Fields）}\label{ux603bux5e8fux4e0eux6709ux5e8fux57dftotal-orders-and-ordered-fields}

\subsubsection{total order}\label{total-order}

一个偏序如果满足``任意 x,y 总有 x≤y 或 y≤x''，就叫全序。\\
整数与有理数在通常顺序下都是全序。

\subsubsection{ordered field}\label{ordered-field}

一个集合若有：

\begin{itemize}
\tightlist
\item
  元素 0,1
\item
  加法与乘法
\item
  域的公理（交换、结合、分配、加法逆元、非零乘法逆元等）
\item
  一个与运算兼容的全序
\end{itemize}

则称为有序域。

Q 是 ordered field。\\
R 也会是 ordered field，但还多出 completeness。

\subsection{4.3
上界、下界、sup、inf}\label{ux4e0aux754cux4e0bux754csupinf}

设 Y⊂X，X 是偏序集。

\begin{itemize}
\tightlist
\item
  upper bound：对所有 y∈Y，都有 y≤x；
\item
  lower bound：对所有 y∈Y，都有 y≥x；
\item
  supremum：所有上界里最小的那个；
\item
  infimum：所有下界里最大的那个。
\end{itemize}

\subsubsection{注意}\label{ux6ce8ux610f}

sup(Y) 不一定在 Y 里面。\\
例如在 \(Q\) 中，正有理数集合 \(Q^+\) 的 inf 是 \(0\)，但 \(0\) 不属于
\(Q^+\)。

\begin{center}\rule{0.5\linewidth}{0.5pt}\end{center}

\subsection{4.4 Lecture Notes 习题详解（Exercises
39--44）}\label{lecture-notes-ux4e60ux9898ux8be6ux89e3exercises-3944}

\subsubsection{Exercise 39}\label{exercise-39}

\textbf{题意}：若 (X,≤) 是全序，则任意子集 Y 继承的限制顺序也是全序。

\textbf{解答}：\\
取任意 \(y_1,y_2\in Y\)。\\
由于 Y⊂X，且 X 是全序，所以在 X 中必有 \(y_1\le y_2\) 或
\(y_2\le y_1\)。\\
而 Y 上的限制顺序就是沿用 X 上的同一比较关系，\\
故在 Y 中也有 \(y_1\le y_2\) 或 \(y_2\le y_1\)。\\
因此 Y 也是全序集。

\subsubsection{Exercise 40}\label{exercise-40}

\textbf{题意}：证明 Proposition 5。

\textbf{内容}： - 若 x≤y，则 x+z≤y+z； - 若 x≥0 且 y≥0，则 xy≥0。

\textbf{解答}：

\begin{enumerate}
\def\labelenumi{\arabic{enumi}.}
\item
  若 x≤y，则 0≤y-x。\\
  两边加上 z 不改变差值： (y+z) - (x+z) = y-x ≥ 0\\
  因此 x+z ≤ y+z。
\item
  若 x≥0 且 y≥0，则依有序域定义中``非负数乘积非负''的公理，得 xy≥0。
\end{enumerate}

\subsubsection{Exercise 41}\label{exercise-41}

\textbf{题意}：证明 Z 没有最大元也没有最小元。

\textbf{解答}：\\
任取 n∈Z。

\begin{itemize}
\tightlist
\item
  n+1 仍是整数，且 n+1\textgreater n，所以 n 不可能是最大元；
\item
  n-1 仍是整数，且 n-1\textless n，所以 n 不可能是最小元。
\end{itemize}

因此 Z 没有 maximum，也没有 minimum。

\subsubsection{Exercise 42}\label{exercise-42}

\textbf{题意}：证明若 s²\textgreater2，则 s 是\\
S=\{x∈Q \textbar{} x²\textless2\} 的上界。

\textbf{解答}：\\
要证：对所有 x∈S，都有 x≤s。

反设存在 x∈S 使 x\textgreater s。\\
由于 x\textgreater s≥0（可只考虑正上界部分；若 s\textless0 则更不可能是
x²\textgreater2 的情形对应的 relevant upper bound），\\
平方函数在正有理数上单调递增，于是 x² \textgreater{} s² \textgreater{} 2
这与 x∈S，即 x²\textless2 矛盾。\\
故对所有 x∈S，x≤s，所以 s 是上界。

\subsubsection{Exercise 43}\label{exercise-43}

\textbf{题意}：若 s²\textless2，证明存在 r∈Q，s\textless r 且
r²\textless2，因此 s 不是上界。

\textbf{解答}：\\
设 M=2-s²\textgreater0。\\
我们要找一个很小的 ε\textgreater0，使 r=s+ε 仍满足 r²\textless2。

展开： r² = (s+ε)² = s² + 2sε + ε²

要使 r²\textless2，只需 2sε + ε² \textless{} 2-s² = M

也就是 ε(2s+ε) \textless{} M

notes 提示我们先选 ε 足够小，使 ε\textless2s；\\
则 2s+ε \textless{} 4s，所以只要再让 ε \textless{} M/(4s) 就能保证
ε(2s+ε) \textless{} M。

由于 Proposition 7 说明任给正有理数都能找到更小的正有理数，\\
所以我们能选到满足 0\textless ε\textless2s 且 ε\textless M/(4s) 的有理数
ε。\\
于是令 r=s+ε，即得 r\textgreater s 且 r²\textless2。\\
因此 s 不是 S 的上界。

\subsubsection{Exercise 44}\label{exercise-44}

\textbf{题意}：若 s²\textgreater2，证明存在 r∈Q，r\textless s 且 r 仍是
S 的上界，因此 s 不是 supremum。

\textbf{解答}：\\
本题与 Exercise 43
对偶：我们想把一个``过大的上界''稍微往左挪一点，但仍保持是上界。

设 r = s - ε 其中 ε\textgreater0 很小。\\
我们希望 r²\textgreater2。展开： r² = (s-ε)² = s² - 2sε + ε²

只要保证 2sε - ε² \textless{} s² - 2 即可推出 r²\textgreater2。\\
因为 s²-2\textgreater0，所以可选足够小的 ε\textgreater0 使上式成立。\\
于是 r²\textgreater2。根据 Exercise 42，任何平方大于 2 的有理数都是 S
的上界，所以 r 仍是上界。\\
但 r\textless s，说明 s 不是最小上界。\\
故 s 不可能是 supremum。

\begin{center}\rule{0.5\linewidth}{0.5pt}\end{center}

\subsection{4.5 为什么 Q 不完备（Why Q is Not
Complete）}\label{ux4e3aux4ec0ux4e48-q-ux4e0dux5b8cux5907why-q-is-not-complete}

考虑 S = \{x∈Q \textbar{} x²\textless2\}

\begin{itemize}
\tightlist
\item
  它有上界：例如任何满足 s²\textgreater2 的有理数都是上界；
\item
  但它在 Q 中没有 supremum。
\end{itemize}

为什么没有？分两种情况：

\begin{enumerate}
\def\labelenumi{\arabic{enumi}.}
\tightlist
\item
  若候选 s 满足 s²\textless2，Exercise 43 告诉你还能往右走一点，且仍留在
  S 中，所以 s 不是上界；
\item
  若候选 s 满足 s²\textgreater2，Exercise 44
  告诉你还能往左挪一点，且仍是上界，所以 s 不是最小上界；
\item
  s²=2 在 Q 中根本不存在。
\end{enumerate}

因此没有有理数能充当 sup(S)。\\
这就是``Q 不完备''的最经典反例。

\subsection{4.6 实数的公理化定义（Axioms for the
Reals）}\label{ux5b9eux6570ux7684ux516cux7406ux5316ux5b9aux4e49axioms-for-the-reals}

实数模型 R 需要满足两点：

\begin{enumerate}
\def\labelenumi{\arabic{enumi}.}
\tightlist
\item
  (R,0,1,+,·,≤) 是一个有序域 (ordered field)；
\item
  (R,≤) 是完备的 (complete)。
\end{enumerate}

这实际上告诉我们：\\
\textbf{实数 = 有理数的有序代数结构 + 不再有``缺口''的顺序结构。}

\subsection{4.7 §4.7
的理解：第一次构造实数的尝试}\label{ux7684ux7406ux89e3ux7b2cux4e00ux6b21ux6784ux9020ux5b9eux6570ux7684ux5c1dux8bd5}

\subsubsection{这一节到底想说什么}\label{ux8fd9ux4e00ux8282ux5230ux5e95ux60f3ux8bf4ux4ec0ux4e48}

它不是让你真正定义 Dedekind cuts，而是告诉你一种动机：

一个实数 r，虽然可能写成无限小数 10.234890234809\ldots{}

但我们永远只能通过``有理数近似''去逼近它：

\begin{itemize}
\tightlist
\item
  10 \textless{} r \textless{} 11
\item
  10.2 \textless{} r \textless{} 10.3
\item
  10.23 \textless{} r \textless{} 10.24
\item
  \ldots{}
\end{itemize}

所以一个实数可以被理解为：\\
\textbf{它把有理数切成``在它左边的一堆''和``在它右边的一堆''。}

也就是说，实数不是某个神秘的``无限小数字串''，\\
而是一个决定了``哪些有理数比它小，哪些比它大''的对象。

\subsubsection{你在期中里要掌握到什么程度}\label{ux4f60ux5728ux671fux4e2dux91ccux8981ux638cux63e1ux5230ux4ec0ux4e48ux7a0bux5ea6}

到 §4.7 为止，你应该会说：

\begin{itemize}
\tightlist
\item
  实数构造的目标是填补 Q 的 gaps；
\item
  小数展开只是动机，不是最严谨定义；
\item
  一个实数可以由它左边 / 右边的有理逼近来刻画；
\item
  Dedekind cuts 是下一节的正式化，但\textbf{不考}。
\end{itemize}

\begin{center}\rule{0.5\linewidth}{0.5pt}\end{center}

\section{第五章
高频证明模板与临考整理}\label{ux7b2cux4e94ux7ae0-ux9ad8ux9891ux8bc1ux660eux6a21ux677fux4e0eux4e34ux8003ux6574ux7406}

\subsection{5.1
如何证明``是等价关系''}\label{ux5982ux4f55ux8bc1ux660eux662fux7b49ux4ef7ux5173ux7cfb}

永远按三步写：

\begin{enumerate}
\def\labelenumi{\arabic{enumi}.}
\tightlist
\item
  Reflexive：xRx
\item
  Symmetric：xRy ⇒ yRx
\item
  Transitive：xRy 且 yRz ⇒ xRz
\end{enumerate}

对 \textasciitilde\_Z、\textasciitilde\_Q、模 m 同余，全都用这个模板。

\subsection{5.2
如何证明``是偏序''}\label{ux5982ux4f55ux8bc1ux660eux662fux504fux5e8f}

永远按三步写：

\begin{enumerate}
\def\labelenumi{\arabic{enumi}.}
\tightlist
\item
  Reflexive
\item
  Antisymmetric
\item
  Transitive
\end{enumerate}

别把 symmetric 和 antisymmetric 混淆。\\
``对称''是双向同真；\\
``反对称''是双向同时成立时必须相等。

\subsection{5.3
如何证明``良定义（well-defined）''}\label{ux5982ux4f55ux8bc1ux660eux826fux5b9aux4e49well-defined}

商集问题必考模板：

\begin{itemize}
\tightlist
\item
  先写``设换代表元后仍等价''；
\item
  再写``要证新定义得到的结果等价''；
\item
  把旧的等价关系式代入；
\item
  通过代数变形推出目标。
\end{itemize}

关键词：\\
\textbf{定义在等价类上的对象，必须与代表元无关。}

\subsection{5.4
如何证明``函数可逆''}\label{ux5982ux4f55ux8bc1ux660eux51fdux6570ux53efux9006}

\begin{itemize}
\tightlist
\item
  先证 injective；
\item
  再证 surjective；
\item
  然后定义逆函数；
\item
  最后验证左右逆。
\end{itemize}

记住： f 可逆 ⇔ f 是 bijection。

\subsection{5.5 归纳法模板}\label{ux5f52ux7eb3ux6cd5ux6a21ux677f}

要证明 ∀n P(n)：

\begin{itemize}
\tightlist
\item
  base case
\item
  induction hypothesis
\item
  induction step
\item
  conclude by induction
\end{itemize}

如果命题含两个变量，如 a+b=b+a，通常对其中一个做归纳，另一个任意。

\subsection{5.6
上确界题的思路}\label{ux4e0aux786eux754cux9898ux7684ux601dux8def}

证明某个 s 是 sup(Y)，通常要两件事都做：

\begin{enumerate}
\def\labelenumi{\arabic{enumi}.}
\tightlist
\item
  \textbf{s 是上界}

  \begin{itemize}
  \tightlist
  \item
    对任意 y∈Y，证明 y≤s
  \end{itemize}
\item
  \textbf{s 是最小上界}

  \begin{itemize}
  \tightlist
  \item
    取任意 t\textless s，证明 t 不是上界\\
    或取任意上界 u，证明 s≤u
  \end{itemize}
\end{enumerate}

若要证明``不存在 supremum in Q''，最常见方法就是：

\begin{itemize}
\tightlist
\item
  若候选太小，说明它根本不是上界；
\item
  若候选太大，说明它不是最小上界；
\item
  若理想值根本不在 Q 中，就说明 gap 出现了。
\end{itemize}

\begin{center}\rule{0.5\linewidth}{0.5pt}\end{center}

\section{第六章
期中前最后一遍必看清单}\label{ux7b2cux516dux7ae0-ux671fux4e2dux524dux6700ux540eux4e00ux904dux5fc5ux770bux6e05ux5355}

\subsection{6.1
你必须能立刻写出的公式}\label{ux4f60ux5fc5ux987bux80fdux7acbux523bux5199ux51faux7684ux516cux5f0f}

\begin{itemize}
\tightlist
\item
  A → B ≡ ¬A ∨ B
\item
  ¬(A ∧ B) ≡ ¬A ∨ ¬B
\item
  ¬(A ∨ B) ≡ ¬A ∧ ¬B
\item
  ¬∀xP(x) ≡ ∃x¬P(x)
\item
  ¬∃xP(x) ≡ ∀x¬P(x)
\item
  {[}(a,b){]} + {[}(c,d){]} = {[}(a+c,b+d){]} （整数）
\item
  {[}(a,b){]} · {[}(c,d){]} = {[}(ac+bd,ad+bc){]} （整数）
\item
  {[}(a,b){]} + {[}(c,d){]} = {[}(ad+bc,bd){]} （有理数）
\item
  {[}(a,b){]} · {[}(c,d){]} = {[}(ac,bd){]} （有理数）
\end{itemize}

\subsection{6.2
你必须能口头解释清楚的概念}\label{ux4f60ux5fc5ux987bux80fdux53e3ux5934ux89e3ux91caux6e05ux695aux7684ux6982ux5ff5}

\begin{itemize}
\tightlist
\item
  proposition 和 open sentence 的区别
\item
  implication 为什么在前件假时为真
\item
  predicate 与 quantifier 是什么
\item
  子集 / 真子集 / 幂集
\item
  image 与 preimage 的区别
\item
  injective / surjective / bijective
\item
  partial order 与 equivalence relation 的区别
\item
  quotient set 是什么
\item
  Peano axioms 的作用
\item
  well-defined 的意义
\item
  completeness 的含义
\item
  为什么 Q 不完备
\end{itemize}

\subsection{6.3
你必须会写的标准证明}\label{ux4f60ux5fc5ux987bux4f1aux5199ux7684ux6807ux51c6ux8bc1ux660e}

\begin{itemize}
\tightlist
\item
  模 m 同余是等价关系
\item
  A⊂B 且 B⊂A 推出 A=B
\item
  0+n=n（归纳）
\item
  加法交换律（归纳）
\item
  gcd(a,b)=gcd(b,c)
\item
  √2 不在 Q
\item
  S=\{q∈Q:q²\textless2\} 在 Q 中没有 supremum
\end{itemize}

\subsection{6.4
最容易错的点}\label{ux6700ux5bb9ux6613ux9519ux7684ux70b9}

\begin{enumerate}
\def\labelenumi{\arabic{enumi}.}
\item
  把 \(A\leftrightarrow B\) 和 \(\varphi\equiv\psi\) 混为一谈\\
  前者是公式，后者是``两个公式逻辑等价''的元层表述。
\item
  忘记 implication 的真值规则\\
  A→B 只有在 A 真且 B 假时为假。
\item
  把 ∀x∃y 与 ∃y∀x 搞混\\
  前者允许 y 随 x 变化，后者不允许。
\item
  把 preimage 当成 inverse function\\
  \(f^{-1}(B)\) 总是有意义；真正的逆函数只在 bijection 时存在。
\item
  商集上定义运算没去证良定义\\
  这是结构构造题的大忌。
\item
  误以为 supremum 一定属于集合本身\\
  错。它只需要属于外部偏序集 X。
\end{enumerate}

\begin{center}\rule{0.5\linewidth}{0.5pt}\end{center}

\section{结语}\label{ux7ed3ux8bed}

如果你把这份讲义真正吃透，那么你对期中前这门课的理解应当已经不是``会算几道题''，而是：

\begin{itemize}
\tightlist
\item
  会把自然语言翻译成逻辑；
\item
  会把代数结构翻译成集合论定义；
\item
  会把``数''理解成由等价关系构造出来的对象；
\item
  会从``顺序结构''角度理解为什么要引入实数。
\end{itemize}

这门课最美的地方，就在于它把你熟悉得不能再熟悉的数字与证明，重新拆开，告诉你：\\
原来一切都可以从最基本的逻辑与集合开始，一步一步严谨地造出来。

\subsection{附注}\label{ux9644ux6ce8}

本讲义严格排除了 Dedekind cuts
作为期中范围之外内容；同时只纳入当前会话中可用的 MATH1090 材料（Lecture
Notes、HW1、HW2、Worksheet
1--5）。若之后补充更多本课程作业，可在此版本上继续扩展。

\section{附录 A 逐节定义与命题速查（Definition-by-Definition Quick
Reference）}\label{ux9644ux5f55-a-ux9010ux8282ux5b9aux4e49ux4e0eux547dux9898ux901fux67e5definition-by-definition-quick-reference}

\subsection{A.1 Logic 速查}\label{a.1-logic-ux901fux67e5}

\subsubsection{命题（proposition）}\label{ux547dux9898proposition-1}

能判定真假的陈述句。

\subsubsection{Boolean formula}\label{boolean-formula}

由命题变量与 ¬, ∧, ∨, →, ↔ 按语法规则拼成的公式。

\subsubsection{tautology}\label{tautology}

无论变量真假如何取值都为真的公式。

\subsubsection{contradiction}\label{contradiction}

无论变量真假如何取值都为假的公式。

\subsubsection{predicate}\label{predicate}

依赖于变量的性质，例如 P(x), R(x,y)。

\subsubsection{universal quantifier / existential
quantifier}\label{universal-quantifier-existential-quantifier}

\begin{itemize}
\tightlist
\item
  ∀x：对所有 x
\item
  ∃x：存在某个 x
\end{itemize}

\subsubsection{exactly one}\label{exactly-one}

∃!x P(x) 表示``存在且仅存在一个 x 使 P(x) 成立''。

\subsection{A.2 Sets 速查}\label{a.2-sets-ux901fux67e5}

\subsubsection{set}\label{set}

集合是对象的汇集。

\subsubsection{subset}\label{subset}

A⊂B 表示 A 中每个元素都属于 B。

\subsubsection{union / intersection /
difference}\label{union-intersection-difference}

\begin{itemize}
\tightlist
\item
  A∪B：至少属于一个
\item
  A∩B：同时属于两者
\item
  \(A\setminus B\)：属于 A 但不属于 B
\end{itemize}

\subsubsection{Cartesian product}\label{cartesian-product}

A×B=\{(a,b)\textbar a∈A,b∈B\}

\subsubsection{power set}\label{power-set}

P(A)=\{X\textbar X⊂A\}

\subsection{A.3 Functions and Relations
速查}\label{a.3-functions-and-relations-ux901fux67e5}

\subsubsection{function}\label{function}

f⊂X×Y，且每个 x∈X 对应唯一 y∈Y。

\subsubsection{image / preimage}\label{image-preimage-1}

\begin{itemize}
\tightlist
\item
  f(A)：A 的像
\item
  f\^{}-¹(B)：B 的原像
\end{itemize}

\subsubsection{injective / surjective /
bijective}\label{injective-surjective-bijective-1}

\begin{itemize}
\tightlist
\item
  injective：不同输入对应不同输出
\item
  surjective：每个目标元素都被命中
\item
  bijective：既单又满
\end{itemize}

\subsubsection{relation}\label{relation}

R⊂X×Y；若 (x,y)∈R，记作 xRy。

\subsubsection{reflexive / symmetric / antisymmetric /
transitive}\label{reflexive-symmetric-antisymmetric-transitive}

\begin{itemize}
\tightlist
\item
  reflexive：xRx
\item
  symmetric：xRy ⇒ yRx
\item
  antisymmetric：xRy 且 yRx ⇒ x=y
\item
  transitive：xRy 且 yRz ⇒ xRz
\end{itemize}

\subsubsection{partial order}\label{partial-order}

自反 + 反对称 + 传递。

\subsubsection{equivalence relation}\label{equivalence-relation}

自反 + 对称 + 传递。

\subsubsection{equivalence class}\label{equivalence-class}

{[}x{]}\_R = \{y \textbar{} xRy\}

\subsubsection{quotient set}\label{quotient-set}

X/\textasciitilde{} = ``所有等价类''的集合。

\subsection{A.4 Naturals, Integers, Rationals
速查}\label{a.4-naturals-integers-rationals-ux901fux67e5}

\subsubsection{model of natural numbers}\label{model-of-natural-numbers}

由 (N,0,S) 组成并满足 Peano axioms。

\subsubsection{principle of induction}\label{principle-of-induction}

若 P(0) 成立且 P(x)⇒P(S(x))，则 ∀xP(x)。

\subsubsection{integers as quotient}\label{integers-as-quotient}

Z = N² / \textasciitilde\_Z，其中\\
(a,b)\textasciitilde\_Z(c,d) ⇔ a+d=b+c

\subsubsection{rationals as quotient}\label{rationals-as-quotient}

Q = (\(Z\times (Z\setminus\{0\})\)) / \textasciitilde\_Q，其中\\
(a,b)\textasciitilde\_Q(c,d) ⇔ ad=bc

\subsubsection{well-defined}\label{well-defined}

在等价类上定义的对象不依赖代表元。

\subsection{A.5 Reals 速查}\label{a.5-reals-ux901fux67e5}

\subsubsection{total order}\label{total-order-1}

任意两个元素都可比较。

\subsubsection{upper bound / lower bound}\label{upper-bound-lower-bound}

对整个子集从上 / 从下控制住的元素。

\subsubsection{supremum / infimum}\label{supremum-infimum}

最小上界 / 最大下界。

\subsubsection{complete}\label{complete}

每个有上界的子集都有 supremum；每个有下界的子集都有 infimum。

\subsubsection{ordered field}\label{ordered-field-1}

既是域，又有与加法和乘法兼容的全序。

\subsubsection{model of the real
numbers}\label{model-of-the-real-numbers}

有序域 + 完备性。

\section{附录 B 60 题快问快答（Rapid-Fire
Self-Test）}\label{ux9644ux5f55-b-60-ux9898ux5febux95eeux5febux7b54rapid-fire-self-test}

\subsection{B.1 Logic}\label{b.1-logic}

\begin{enumerate}
\def\labelenumi{\arabic{enumi}.}
\item
  什么是 proposition？\\
  能判定真假值的陈述句。
\item
  ``x+1=3'' 在未指定 x 时是 proposition 吗？\\
  不是，通常是开放句。
\item
  A→B 什么时候是假？\\
  仅当 A 真且 B 假。
\item
  A→B 与什么等价？\\
  ¬A∨B。
\item
  ¬(A→B) 与什么等价？\\
  A∧¬B。
\item
  ¬(A∨B) 与什么等价？\\
  ¬A∧¬B。
\item
  ¬(A∧B) 与什么等价？\\
  ¬A∨¬B。
\item
  A↔B 表示什么？\\
  A 与 B 同真同假。
\item
  tautology 是什么？\\
  恒真式。
\item
  contradiction 是什么？\\
  恒假式。
\item
  modus tollens 的形式是什么？\\
  B→A，¬A，因此 ¬B。
\item
  affirming the consequent 为什么错？\\
  因为后件可能由别的原因导致。
\item
  ∀xP(x) 的否定是什么？\\
  ∃x¬P(x)。
\item
  ∃xP(x) 的否定是什么？\\
  ∀x¬P(x)。
\item
  ∀x∃yP(x,y) 与 ∃y∀xP(x,y) 一般等价吗？\\
  一般不等价。
\end{enumerate}

\subsection{B.2 Sets and Functions}\label{b.2-sets-and-functions}

\begin{enumerate}
\def\labelenumi{\arabic{enumi}.}
\setcounter{enumi}{15}
\item
  证明两个集合相等最常用的方法是什么？\\
  双向包含或元素追踪。
\item
  P(A) 表示什么？\\
  A 的幂集。
\item
  若 \textbar A\textbar=n，则 \textbar P(A)\textbar=？\\
  2\^{}n。
\item
  A×B 的元素是什么？\\
  有序对 (a,b)。
\item
  函数在集合论里是什么？\\
  X×Y 的特殊子集。
\item
  什么叫 injective？\\
  不同输入给不同输出。
\item
  什么叫 surjective？\\
  目标集合每个元素都是输出。
\item
  可逆函数一定是什么？\\
  bijection。
\item
  preimage 是否一定存在？\\
  对任意函数和任意目标子集都存在。
\item
  inverse function 是否一定存在？\\
  不一定，只有 bijection 才有。
\item
  左逆存在能推出什么？\\
  函数是单射。
\item
  右逆存在能推出什么？\\
  函数是满射。
\end{enumerate}

\subsection{B.3 Relations}\label{b.3-relations}

\begin{enumerate}
\def\labelenumi{\arabic{enumi}.}
\setcounter{enumi}{27}
\item
  relation 是什么？\\
  某个乘积集合的子集。
\item
  partial order 三个条件是什么？\\
  reflexive, antisymmetric, transitive。
\item
  equivalence relation 三个条件是什么？\\
  reflexive, symmetric, transitive。
\item
  ≤ 是 symmetric 吗？\\
  不是。
\item
  ≤ 是 antisymmetric 吗？\\
  是。
\item
  模 m 同余是哪种关系？\\
  equivalence relation。
\item
  等价类有什么作用？\\
  把集合分块，并形成 quotient set。
\item
  Hasse diagram 画的是什么？\\
  偏序结构中的覆盖关系。
\end{enumerate}

\subsection{B.4 Naturals and
Induction}\label{b.4-naturals-and-induction}

\begin{enumerate}
\def\labelenumi{\arabic{enumi}.}
\setcounter{enumi}{35}
\item
  Peano 公理中最关键的一条是什么？\\
  principle of induction。
\item
  加法递归定义是什么？\\
  a+0=a，a+S(b)=S(a+b)。
\item
  乘法递归定义是什么？\\
  a·0=0，a·S(b)=a+(a·b)。
\item
  证明 0+n=n 用什么方法？\\
  对 n 归纳。
\item
  证明 a+b=b+a 通常对谁归纳？\\
  通常对 b 归纳。
\end{enumerate}

\subsection{B.5 Integers and
Rationals}\label{b.5-integers-and-rationals}

\begin{enumerate}
\def\labelenumi{\arabic{enumi}.}
\setcounter{enumi}{40}
\item
  Z 是如何构造的？\\
  作为 N² 上等价关系的商集。
\item
  (a,b) 在 Z 中直观表示什么？\\
  a-b。
\item
  Z 上的等价关系是什么？\\
  (a,b)\textasciitilde(c,d) ⇔ a+d=b+c。
\item
  Q 是如何构造的？\\
  作为 \(Z\times (Z\setminus\{0\})\) 上等价关系的商集。
\item
  (a,b) 在 Q 中直观表示什么？\\
  a/b。
\item
  Q 上的等价关系是什么？\\
  (a,b)\textasciitilde(c,d) ⇔ ad=bc。
\item
  为什么要证明运算 well-defined？\\
  因为同一个等价类有多个代表元。
\item
  非零有理数为何有乘法逆元？\\
  因为 a/b 的逆是 b/a（当 a≠0）。
\item
  证明 √2 不在 Q 的关键矛盾是什么？\\
  最简分数的分子分母都会变成偶数。
\end{enumerate}

\subsection{B.6 Reals and
Completeness}\label{b.6-reals-and-completeness}

\begin{enumerate}
\def\labelenumi{\arabic{enumi}.}
\setcounter{enumi}{49}
\item
  total order 比 partial order 多了什么？\\
  任意两元素都可比较。
\item
  upper bound 与 maximum 一样吗？\\
  不一样；upper bound 可以不在原集合里。
\item
  supremum 一定在集合中吗？\\
  不一定。
\item
  infimum 一定在集合中吗？\\
  不一定。
\item
  \(Q^+\) 的 inf 是多少？\\
  0。
\item
  Z 有 maximum 吗？\\
  没有。
\item
  Z 有 minimum 吗？\\
  没有。
\item
  complete 的定义是什么？\\
  有上界就有 sup，有下界就有 inf。
\item
  Q 为什么不 complete？\\
  因为集合 \{q∈Q\textbar q²\textless2\} 在 Q 中无 supremum。
\item
  R 的公理化定义是什么？\\
  ordered field + completeness。
\item
  §4.7 想表达的核心观念是什么？\\
  实数可由不断细化的有理逼近来刻画，它把 Q 分成左边与右边两部分。
\end{enumerate}

\section{附录 C 考前半小时速记版（Ultra-Short Last-Minute
Sheet）}\label{ux9644ux5f55-c-ux8003ux524dux534aux5c0fux65f6ux901fux8bb0ux7248ultra-short-last-minute-sheet}

\subsection{C.1
五句最重要的话}\label{c.1-ux4e94ux53e5ux6700ux91cdux8981ux7684ux8bdd}

\begin{enumerate}
\def\labelenumi{\arabic{enumi}.}
\tightlist
\item
  implication：只有``前真后假''时假。\\
\item
  量词否定：¬∀=∃¬，¬∃=∀¬。\\
\item
  函数可逆当且仅当 bijection。\\
\item
  商集上定义运算必须先证 well-defined。\\
\item
  R 比 Q 多出来的关键是 completeness。
\end{enumerate}

\subsection{C.2
五类最重要题型}\label{c.2-ux4e94ux7c7bux6700ux91cdux8981ux9898ux578b}

\begin{enumerate}
\def\labelenumi{\arabic{enumi}.}
\tightlist
\item
  真值表与逻辑等价\\
\item
  自然语言与谓词公式互译\\
\item
  等价关系 / 偏序关系判定\\
\item
  归纳法证明\\
\item
  ``Q 不完备''与 √2 非有理
\end{enumerate}

\subsection{C.3
五个最危险混淆点}\label{c.3-ux4e94ux4e2aux6700ux5371ux9669ux6df7ux6dc6ux70b9}

\begin{enumerate}
\def\labelenumi{\arabic{enumi}.}
\tightlist
\item
  \(A\leftrightarrow B\) 与 \(\varphi\equiv\psi\)\\
\item
  preimage 与 inverse\\
\item
  symmetric 与 antisymmetric\\
\item
  ∀x∃y 与 ∃y∀x\\
\item
  upper bound / maximum / supremum
\end{enumerate}

\section{附录 D 补充严解：把 lecture notes
中``只给提示''或``略写''的部分完整写开}\label{ux9644ux5f55-d-ux8865ux5145ux4e25ux89e3ux628a-lecture-notes-ux4e2dux53eaux7ed9ux63d0ux793aux6216ux7565ux5199ux7684ux90e8ux5206ux5b8cux6574ux5199ux5f00}

\subsection{D.1 命题逻辑等价式的真值表全验算（对应 Lecture Exercise
2）}\label{d.1-ux547dux9898ux903bux8f91ux7b49ux4ef7ux5f0fux7684ux771fux503cux8868ux5168ux9a8cux7b97ux5bf9ux5e94-lecture-exercise-2}

下面把期中前最常用的几条逻辑等价式完整验算一次。考试时你未必需要每次都从头列，但至少应会在纸上重建这些表。

\subsubsection{D.1.1 蕴含改写：A→B ≡
¬A∨B}\label{d.1.1-ux8574ux542bux6539ux5199ab-ab}

\begin{Shaded}
\begin{Highlighting}[]
\NormalTok{A  B | ¬A | ¬A∨B | A→B}
\NormalTok{F  F |  T |   T  |  T}
\NormalTok{F  T |  T |   T  |  T}
\NormalTok{T  F |  F |   F  |  F}
\NormalTok{T  T |  F |   T  |  T}
\end{Highlighting}
\end{Shaded}

最后两列完全相同，所以\\
A→B ≡ ¬A∨B。

\subsubsection{D.1.2 德摩根律：¬(A∨B) ≡
¬A∧¬B}\label{d.1.2-ux5fb7ux6469ux6839ux5f8bab-ab}

\begin{Shaded}
\begin{Highlighting}[]
\NormalTok{A  B | A∨B | ¬(A∨B) | ¬A | ¬B | ¬A∧¬B}
\NormalTok{F  F |  F  |    T   |  T |  T |   T}
\NormalTok{F  T |  T  |    F   |  T |  F |   F}
\NormalTok{T  F |  T  |    F   |  F |  T |   F}
\NormalTok{T  T |  T  |    F   |  F |  F |   F}
\end{Highlighting}
\end{Shaded}

故\\
¬(A∨B) ≡ ¬A∧¬B。

\subsubsection{D.1.3 德摩根律：¬(A∧B) ≡
¬A∨¬B}\label{d.1.3-ux5fb7ux6469ux6839ux5f8bab-ab}

\begin{Shaded}
\begin{Highlighting}[]
\NormalTok{A  B | A∧B | ¬(A∧B) | ¬A | ¬B | ¬A∨¬B}
\NormalTok{F  F |  F  |    T   |  T |  T |   T}
\NormalTok{F  T |  F  |    T   |  T |  F |   T}
\NormalTok{T  F |  F  |    T   |  F |  T |   T}
\NormalTok{T  T |  T  |    F   |  F |  F |   F}
\end{Highlighting}
\end{Shaded}

故\\
¬(A∧B) ≡ ¬A∨¬B。

\subsubsection{D.1.4 逆否命题：A→B ≡
¬B→¬A}\label{d.1.4-ux9006ux5426ux547dux9898ab-ba}

\begin{Shaded}
\begin{Highlighting}[]
\NormalTok{A  B | A→B | ¬B | ¬A | ¬B→¬A}
\NormalTok{F  F |  T  |  T |  T |   T}
\NormalTok{F  T |  T  |  F |  T |   T}
\NormalTok{T  F |  F  |  T |  F |   F}
\NormalTok{T  T |  T  |  F |  F |   T}
\end{Highlighting}
\end{Shaded}

故\\
A→B ≡ ¬B→¬A。

\subsubsection{D.1.5
交换律与结合律为何``肉眼可见''地成立}\label{d.1.5-ux4ea4ux6362ux5f8bux4e0eux7ed3ux5408ux5f8bux4e3aux4f55ux8089ux773cux53efux89c1ux5730ux6210ux7acb}

\begin{itemize}
\tightlist
\item
  A∧B 与 B∧A：因为``同时为真''不依赖先后顺序；
\item
  A∨B 与 B∨A：因为``至少一个为真''不依赖先后顺序；
\item
  A∧(B∧C) 与 (A∧B)∧C：都等价于``三个都真''；
\item
  A∨(B∨C) 与 (A∨B)∨C：都等价于``三个里至少一个真''。
\end{itemize}

\subsubsection{D.1.6
分配律的语义解释}\label{d.1.6-ux5206ux914dux5f8bux7684ux8bedux4e49ux89e3ux91ca}

A∧(B∨C)\\
表示：A 真，而且 B、C 至少一个真。\\
这与 (A∧B)∨(A∧C) 完全一致：要么 A 和 B 真，要么 A 和 C 真。

同理\\
A∨(B∧C) ≡ (A∨B)∧(A∨C)。

\begin{center}\rule{0.5\linewidth}{0.5pt}\end{center}

\subsection{D.2 Lecture Exercise 29
的完整证明骨架（自然数上的交换律、结合律）}\label{d.2-lecture-exercise-29-ux7684ux5b8cux6574ux8bc1ux660eux9aa8ux67b6ux81eaux7136ux6570ux4e0aux7684ux4ea4ux6362ux5f8bux7ed3ux5408ux5f8b}

这一题在 notes 里是一句带过，但真正写下来其实是整章最核心的训练之一。

\subsubsection{D.2.1 先备一个引理：S(a)+b =
S(a+b)}\label{d.2.1-ux5148ux5907ux4e00ux4e2aux5f15ux7406sab-sab}

这个引理在 Worksheet 4 Exercise 1 已证明，这里重写一次：

对 b 归纳。

\begin{itemize}
\item
  Base case： S(a)+0 = S(a) = S(a+0)
\item
  Induction step： 假设 S(a)+b = S(a+b)，则 S(a)+S(b) = S(S(a)+b) =
  S(S(a+b)) = S(a+S(b))
\end{itemize}

故对所有 b，S(a)+b = S(a+b)。

\subsubsection{D.2.2
加法交换律：a+b=b+a}\label{d.2.2-ux52a0ux6cd5ux4ea4ux6362ux5f8babba}

对 b 归纳。

\begin{itemize}
\item
  Base case： a+0 = a = 0+a
\item
  Induction hypothesis： 假设对所有 a，a+b=b+a。
\item
  Induction step： a+S(b) = S(a+b) （由加法定义） = S(b+a)
  （由归纳假设） = S(b)+a （由引理）
\end{itemize}

故 a+S(b)=S(b)+a。\\
由归纳法，加法交换律成立。

\subsubsection{D.2.3
加法结合律：a+(b+c)=(a+b)+c}\label{d.2.3-ux52a0ux6cd5ux7ed3ux5408ux5f8babcabc}

对 c 归纳。

\begin{itemize}
\item
  Base case： a+(b+0)=a+b=(a+b)+0
\item
  Induction step： 假设 a+(b+c)=(a+b)+c。 则 a+(b+S(c)) = a+S(b+c) =
  S(a+(b+c)) = S((a+b)+c) = (a+b)+S(c)
\end{itemize}

故成立。

\subsubsection{D.2.4
乘法交换律：a·b=b·a}\label{d.2.4-ux4e58ux6cd5ux4ea4ux6362ux5f8babba}

证明比加法长，因此最稳妥的写法是对 b 归纳。

\begin{itemize}
\item
  Base case： a·0=0，而 0·a=0（可再对 a 归纳证明）。\\
  所以 a·0=0·a。
\item
  Induction step： 假设 a·b=b·a。\\
  则 a·S(b) = a + a·b = a + b·a 我们希望把它改写成 S(b)·a。\\
  由于 S(b)·a 需要对 a 再用归纳或使用已知引理： S(b)·a = b·a + a 于是 a
  + b·a = b·a + a = S(b)·a
\end{itemize}

故 a·S(b)=S(b)·a。

\subsubsection{D.2.5
乘法结合律：a·(b·c)=(a·b)·c}\label{d.2.5-ux4e58ux6cd5ux7ed3ux5408ux5f8babcabc}

对 c 归纳。

\begin{itemize}
\item
  Base case： a·(b·0)=a·0=0=(a·b)·0
\item
  Induction step： 假设 a·(b·c)=(a·b)·c。 则 a·(b·S(c)) = a·(b + b·c) =
  a·b + a·(b·c) （这里用分配律；若尚未证明，也可把分配律先作为 Exercise
  30 完成） = a·b + (a·b)·c = (a·b)·S(c)
\end{itemize}

因此成立。

\textbf{说明}：严格的公理化发展里，Exercise 29 与 Exercise
30（分配律）之间会互相配合，所以在书写顺序上允许先证一个更容易的辅助版本，再回头补齐另一个。考试若不要求到完全形式化级别，把``归纳变量
+ 每一步所用定义''写清通常就够了。

\begin{center}\rule{0.5\linewidth}{0.5pt}\end{center}

\subsection{D.3 Worksheet 3 Exercise 5：N×N 注入 N
的显式证明}\label{d.3-worksheet-3-exercise-5nn-ux6ce8ux5165-n-ux7684ux663eux5f0fux8bc1ux660e}

定义 f(m,n)=((m+n)(m+n+1))/2 + n

这叫 Cantor pairing function。

\subsubsection{为什么它是注入}\label{ux4e3aux4ec0ux4e48ux5b83ux662fux6ce8ux5165}

记 s=m+n。\\
那么所有满足 m+n=s 的点会被排在第 s 条``对角线''上。\\
每条对角线的起点编号是 0 + 1 + 2 + \ldots{} + s = s(s+1)/2
再加上当前在这条对角线上的偏移量 n，就得到 f(m,n)=s(s+1)/2+n。

若 f(m,n)=f(m',n') 则它们必落在同一条对角线（否则前缀和不同），所以 m+n
= m'+n' 再由同一个公式中的偏移量相等得到 n=n' 从而 m=m'。\\
所以 (m,n)=(m',n')，因此 f 是单射。

\begin{center}\rule{0.5\linewidth}{0.5pt}\end{center}

\subsection{D.4 Lecture Exercise 34 与 37：well-defined
的完整书写模板}\label{d.4-lecture-exercise-34-ux4e0e-37well-defined-ux7684ux5b8cux6574ux4e66ux5199ux6a21ux677f}

\subsubsection{D.4.1 整数乘法良定义（Exercise
34）}\label{d.4.1-ux6574ux6570ux4e58ux6cd5ux826fux5b9aux4e49exercise-34}

定义： {[}(a,b){]}·{[}(c,d){]} := {[}(ac+bd, ad+bc){]}

要证：若 (a,b)\textasciitilde\_Z(a',b')，(c,d)\textasciitilde\_Z(c',d')
即 a+b' = b+a'， c+d' = d+c'， 则 (ac+bd, ad+bc) \textasciitilde\_Z
(a'c'+b'd', a'd'+b'c')。

根据 \textasciitilde\_Z 的定义，我们需证明 (ac+bd) + (a'd'+b'c') =
(ad+bc) + (a'c'+b'd')

把左边展开： ac + bd + a'd' + b'c'

把右边展开： ad + bc + a'c' + b'd'

现在利用已知 a+b' = b+a' \ldots\ldots(1) c+d' = d+c' \ldots\ldots(2)

把 (1)(2) 两式相乘： (a+b')(c+d') = (b+a')(d+c')

展开得 ac + ad' + b'c + b'd' = bd + bc' + a'd + a'c'

再通过交换律、结合律重排，可把式子改写为 ac + a'd' + bd + b'c' = ad +
a'c' + bc + b'd'

这正是所需等式的同型展开。\\
因此整数乘法良定义。

\subsubsection{D.4.2 有理数乘法良定义（Exercise
37）}\label{d.4.2-ux6709ux7406ux6570ux4e58ux6cd5ux826fux5b9aux4e49exercise-37}

定义： {[}(a,b){]}·{[}(c,d){]} := {[}(ac, bd){]}

若 (a,b)\textasciitilde\_Q(a',b')，(c,d)\textasciitilde\_Q(c',d') 即 ab'
= ba'，cd' = dc'

则要证 (ac,bd) \textasciitilde\_Q (a'c', b'd') 根据 \textasciitilde\_Q
定义，只需证 (ac)(b'd') = (bd)(a'c')

左边 (ac)(b'd') = (ab')(cd') 右边 (bd)(a'c') = (ba')(dc')

又由于 ab'=ba'，cd'=dc' 可得左右相等。\\
故良定义。

\subsubsection{D.4.3
这类题为什么总是这样写}\label{d.4.3-ux8fd9ux7c7bux9898ux4e3aux4ec0ux4e48ux603bux662fux8fd9ux6837ux5199}

因为``商集上的一个元素''不是某个具体代表元，而是一整个等价类。\\
若你用不同代表元算出不同结果，则这个运算就根本没有定义成功。\\
所以：

\begin{itemize}
\tightlist
\item
  先换代表元；
\item
  再证结果不变；
\end{itemize}

这是 quotient construction 的灵魂。

\begin{center}\rule{0.5\linewidth}{0.5pt}\end{center}

\subsection{D.5 Lecture Exercise 35 与
38：整数与有理数运算性质的更严格展开}\label{d.5-lecture-exercise-35-ux4e0e-38ux6574ux6570ux4e0eux6709ux7406ux6570ux8fd0ux7b97ux6027ux8d28ux7684ux66f4ux4e25ux683cux5c55ux5f00}

\subsubsection{D.5.1
整数加法交换律}\label{d.5.1-ux6574ux6570ux52a0ux6cd5ux4ea4ux6362ux5f8b}

{[}(a,b){]} + {[}(c,d){]} = {[}(a+c, b+d){]} = {[}(c+a, d+b){]} =
{[}(c,d){]} + {[}(a,b){]}

利用的是 N 上加法交换律。

\subsubsection{D.5.2
整数加法结合律}\label{d.5.2-ux6574ux6570ux52a0ux6cd5ux7ed3ux5408ux5f8b}

({[}(a,b){]} + {[}(c,d){]}) + {[}(e,f){]} = {[}(a+c,b+d){]} +
{[}(e,f){]} = {[}(a+c+e, b+d+f){]}

而 {[}(a,b){]} + ({[}(c,d){]} + {[}(e,f){]}) = {[}(a,b){]} +
{[}(c+e,d+f){]} = {[}(a+c+e, b+d+f){]}

所以相等。

\subsubsection{D.5.3
整数乘法交换律}\label{d.5.3-ux6574ux6570ux4e58ux6cd5ux4ea4ux6362ux5f8b}

{[}(a,b){]}·{[}(c,d){]} = {[}(ac+bd, ad+bc){]} = {[}(ca+db, cb+da){]} =
{[}(c,d){]}·{[}(a,b){]}

\subsubsection{D.5.4
整数分配律}\label{d.5.4-ux6574ux6570ux5206ux914dux5f8b}

左边： {[}(a,b){]}·( {[}(c,d){]} + {[}(e,f){]} ) =
{[}(a,b){]}·{[}(c+e,d+f){]} = {[}(a(c+e)+b(d+f), a(d+f)+b(c+e)){]} =
{[}(ac+ae+bd+bf, ad+af+bc+be){]}

右边： {[}(a,b){]}·{[}(c,d){]} + {[}(a,b){]}·{[}(e,f){]} =
{[}(ac+bd,ad+bc){]} + {[}(ae+bf,af+be){]} = {[}(ac+bd+ae+bf,
ad+bc+af+be){]}

两边只差交换结合的重排，因此相等。

\subsubsection{D.5.5
有理数加法交换律}\label{d.5.5-ux6709ux7406ux6570ux52a0ux6cd5ux4ea4ux6362ux5f8b}

{[}(a,b){]} + {[}(c,d){]} = {[}(ad+bc, bd){]} = {[}(cb+da, db){]} =
{[}(c,d){]} + {[}(a,b){]}

\subsubsection{D.5.6
有理数加法结合律}\label{d.5.6-ux6709ux7406ux6570ux52a0ux6cd5ux7ed3ux5408ux5f8b}

左边： ({[}(a,b){]} + {[}(c,d){]}) + {[}(e,f){]} = {[}(ad+bc,bd){]} +
{[}(e,f){]} = {[}((ad+bc)f + (bd)e, bdf){]} = {[}(adf + bcf + bde,
bdf){]}

右边： {[}(a,b){]} + ({[}(c,d){]} + {[}(e,f){]}) = {[}(a,b){]} +
{[}(cf+de, df){]} = {[}(a(df) + b(cf+de), bdf){]} = {[}(adf + bcf + bde,
bdf){]}

故成立。

\subsubsection{D.5.7
有理数乘法结合律与分配律}\label{d.5.7-ux6709ux7406ux6570ux4e58ux6cd5ux7ed3ux5408ux5f8bux4e0eux5206ux914dux5f8b}

完全同理，直接展开即可。\\
例如分配律：

{[}(a,b){]}·( {[}(c,d){]} + {[}(e,f){]} ) = {[}(a,b){]}·{[}(cf+de,
df){]} = {[}(a(cf+de), bdf){]} = {[}(acf+ade, bdf){]}

而 {[}(a,b){]}·{[}(c,d){]} + {[}(a,b){]}·{[}(e,f){]} = {[}(ac,bd){]} +
{[}(ae,bf){]} = {[}(ac·bf + ae·bd, bdbf){]} = {[}(abcf + abde, b²df){]}

在 Q
中分子分母可同时乘上同一个非零整数，因此两式代表同一有理数；更直接的儿童算术直觉是：
a/b · (c/d + e/f) = a/b · ((cf+de)/df) = a(cf+de)/(bdf) = acf/(bdf) +
ade/(bdf) = ac/(bd) + ae/(bf)

\begin{center}\rule{0.5\linewidth}{0.5pt}\end{center}

\subsection{D.6 Lecture Exercise 43 与 44：Q
不完备反例的严密写法}\label{d.6-lecture-exercise-43-ux4e0e-44q-ux4e0dux5b8cux5907ux53cdux4f8bux7684ux4e25ux5bc6ux5199ux6cd5}

令 S = \{x∈Q \textbar{} x² \textless{} 2\}。

\subsubsection{D.6.1 若 s²\textless2，则 s
不是上界}\label{d.6.1-ux82e5-suxb22ux5219-s-ux4e0dux662fux4e0aux754c}

设 M=2-s²\textgreater0。\\
要找有理 ε\textgreater0 使 (s+ε)² \textless{} 2。

展开： (s+ε)² = s² + 2sε + ε²

因此只需 2sε + ε² \textless{} M。

如果先要求 ε \textless{} 1，则 ε² \textless{} ε，故只需再让 (2s+1)ε
\textless{} M。 由于 M/(2s+1) \textgreater{}
0，而有理数在正数之间稠密，可以选到 0 \textless{} ε \textless{} min(1,
M/(2s+1))。 于是 2sε + ε² \textless{} 2sε + ε \textless{} (2s+1)ε
\textless{} M， 所以 (s+ε)² \textless{} s² + M = 2。

令 r=s+ε，则 r\textgreater s 且 r²\textless2，即 r∈S。\\
所以 s 不是上界。

\subsubsection{D.6.2 若 s²\textgreater2，则 s
不是最小上界}\label{d.6.2-ux82e5-suxb22ux5219-s-ux4e0dux662fux6700ux5c0fux4e0aux754c}

设 N=s²-2\textgreater0。\\
要找 ε\textgreater0 很小，使 r=s-ε 仍满足 r²\textgreater2。

展开： (s-ε)² = s² - 2sε + ε²

只需保证 2sε - ε² \textless{} N。

若取 ε\textless1，则 ε²\textgreater0，所以只要 2sε \textless{} N
即可。\\
选 0 \textless{} ε \textless{} min(1, N/(2s)) 便有 (s-ε)² \textgreater{}
s² - N = 2。

于是 r=s-ε 仍满足 r²\textgreater2。\\
根据 Exercise 42，r 仍是 S 的上界。\\
但 r\textless s，所以 s 不是最小上界，即不是 supremum。

\subsubsection{D.6.3 为什么这三题合起来就推出``无
supremum''}\label{d.6.3-ux4e3aux4ec0ux4e48ux8fd9ux4e09ux9898ux5408ux8d77ux6765ux5c31ux63a8ux51faux65e0-supremum}

\begin{itemize}
\tightlist
\item
  平方小于 2 的有理数都不是上界；
\item
  平方大于 2 的有理数都不是最小上界；
\item
  平方等于 2 的有理数不存在。
\end{itemize}

于是 Q 中没有任何元素能当 sup(S)。

\begin{center}\rule{0.5\linewidth}{0.5pt}\end{center}

\subsection{D.7 Worksheet 5 Exercise 4：把``有理区间里充满
gaps''写成一个完整构造}\label{d.7-worksheet-5-exercise-4ux628aux6709ux7406ux533aux95f4ux91ccux5145ux6ee1-gapsux5199ux6210ux4e00ux4e2aux5b8cux6574ux6784ux9020}

题目：任给 a\textless b，证明存在\\
X ⊂ \{x∈Q \textbar{} a\textless x\textless b\} 使 sup(X) 在 Q 中不存在。

\subsubsection{做法：把经典集合 S=\{q∈Q \textbar{} q²\textless2\}
仿射缩放到区间 (a,b)
里}\label{ux505aux6cd5ux628aux7ecfux5178ux96c6ux5408-sqq-quxb22-ux4effux5c04ux7f29ux653eux5230ux533aux95f4-ab-ux91cc}

先取 α = (a+b)/2， β = (b-a)/4。 则 β\textgreater0，且区间 (α-β√2,
α+β√2) 严格包含在 (a,b) 内，因为 β√2 \textless{} (b-a)/2。

定义 X = \{ α + βq \textbar{} q∈Q, q² \textless{} 2 \}。

因为 q²\textless2 时有 -√2 \textless{} q \textless{} √2，\\
所以 α - β√2 \textless{} α + βq \textless{} α + β√2， 故 X⊂(a,b)∩Q。

现在观察：若 X 在 Q 中有 supremum，则记为 s。\\
把它反向缩放回去： T = \{ (x-α)/β \textbar{} x∈X \} = \{ q∈Q \textbar{}
q²\textless2 \} = S 则 s 对应地给出 S 的一个有理 supremum，矛盾。\\
因此 X 在 Q 中没有 supremum。

\subsubsection{这个构造的思想}\label{ux8fd9ux4e2aux6784ux9020ux7684ux601dux60f3}

不是``硬想一个神秘集合''，而是把已知有 gap 的集合 S
通过线性变换搬到任意有理区间里。\\
这就是``有理数到处有 gaps''的真正含义。

\begin{center}\rule{0.5\linewidth}{0.5pt}\end{center}

\subsection{D.8 对开放题与 notes
没有唯一答案的题，应该怎么写}\label{d.8-ux5bf9ux5f00ux653eux9898ux4e0e-notes-ux6ca1ux6709ux552fux4e00ux7b54ux6848ux7684ux9898ux5e94ux8be5ux600eux4e48ux5199}

\subsubsection{Exercise 1（Boolean formulas
自编）}\label{exercise-1boolean-formulas-ux81eaux7f16}

关键不是唯一答案，而是：

\begin{itemize}
\tightlist
\item
  你写的句子要能翻译成符号；
\item
  符号要符合逻辑语法；
\item
  若题目背景允许多个解释，你要把自己的解释说清。
\end{itemize}

\subsubsection{Exercise 21 /
23（关系举例）}\label{exercise-21-23ux5173ux7cfbux4e3eux4f8b}

这类题最好的答法不是只给一个例子，而是按性质分类：

\begin{itemize}
\tightlist
\item
  一个满足该性质的例子；
\item
  一个不满足该性质的例子；
\item
  并明确指出失败在哪里。
\end{itemize}

\subsubsection{Worksheet 3 Exercise 2（Venn
图）}\label{worksheet-3-exercise-2venn-ux56fe}

考试时如果不能画图，也可以用文字写出``谁在谁里面、哪些交为空、哪些两两相交但三重交为空''。\\
只要逻辑关系清楚，文字版一样可以拿分。

\begin{center}\rule{0.5\linewidth}{0.5pt}\end{center}

\subsection{D.9
真正的复习重点：哪些题最值得二刷、三刷}\label{d.9-ux771fux6b63ux7684ux590dux4e60ux91cdux70b9ux54eaux4e9bux9898ux6700ux503cux5f97ux4e8cux5237ux4e09ux5237}

\subsubsection{第一梯队（必须熟到能默写）}\label{ux7b2cux4e00ux68afux961fux5fc5ux987bux719fux5230ux80fdux9ed8ux5199}

\begin{itemize}
\tightlist
\item
  HW1 第 3--6 题
\item
  HW2 全部
\item
  Lecture Exercise 24, 27, 32, 36, 42--44
\item
  Worksheet 4 Exercise 1--2
\item
  Worksheet 5 Exercise 2--4
\end{itemize}

\subsubsection{第二梯队（需要理解结构）}\label{ux7b2cux4e8cux68afux961fux9700ux8981ux7406ux89e3ux7ed3ux6784}

\begin{itemize}
\tightlist
\item
  Lecture Exercise 15, 18, 19
\item
  Lecture Exercise 29--30, 34--35, 37--38
\item
  Worksheet 3 Exercise 3--5
\end{itemize}

\subsubsection{第三梯队（开阔理解、建立直觉）}\label{ux7b2cux4e09ux68afux961fux5f00ux9614ux7406ux89e3ux5efaux7acbux76f4ux89c9}

\begin{itemize}
\tightlist
\item
  Exercise 1, 8, 10, 21, 23
\item
  Worksheet 3 Exercise 2
\end{itemize}

\end{document}
